
\chapter{Limits and Colimits}

\section{Definition}

Given two categories \(\mathcal{I}\) and \(\mathcal{C}\) and
\(v \in \obj{\mathcal{C}}\), we define the functor
\(k_v : \mathcal{I} \to \mathcal{C}\) as the one sending all objects of
\(\mathcal{I}\) to \(v\) and all morphisms of \(\mathcal{C}\) to \(\id_v\).

\begin{definition}
  Let \(X : \mathcal{I} \to \mathcal{C}\) be a functor. A {\em cone} over
  \(X\) is a natural transformation \(k_v \Rightarrow X\). Here,
  \(v \in \obj{\mathcal{C}}\) is called {\em vertex} of the cone. Dually, a {\em
    cocone} under \(X\) is a natural transformation
  \(X \Rightarrow k_v\). Also in this case, we will refer to \(v\) as the vertex
  of the cocone.
\end{definition}

More explicitly, a cone \(\eta : k_v \Rightarrow X\) is a family
\(\set{ \eta_i : v \to X(i) | i \in \obj{\mathcal{I}} }\) such that the following diagram
commutes for every \(f : i \to j\) of \(\mathcal{I}\)
\[
  \begin{tikzcd}[row sep=tiny]
    & X(i) \ar["{X(f)}", dd] \\
    v \ar["{\eta_i}", ur] \ar["{\eta_j}", dr, swap] \\
    & X(j)
  \end{tikzcd}
\]

\begin{definition}[Limits \&
  colimits]\label{definition:LimitsAndColimits}
  A {\em limit} of a functor \(X : \mathcal{I} \to \mathcal{C}\) is some object
  \(v\) of \(\mathcal{C}\) together with a cone \(\lambda : k_v \tto X\) such that:
  % \begin{quotation}
  for any object \(a\) of \(\mathcal{C}\) and \(\mu : k_a \tto X\) there is one and
  only one \(f : a \to v\) of \(\mathcal{C}\) such that
  \[
    \begin{tikzcd}[row sep=tiny]
      a \ar["{\mu_i}", dr] \ar["f", dd, swap] \\
      & X(i) \\
      v \ar["{\lambda_i}", ur, swap]
    \end{tikzcd}
  \]
  commutes for every \(i\) in \(\mathcal{I}\).
  % \end{quotation}
  A {\em colimit}, instead, is an object \(u\) of \(\mathcal{C}\) together with a
  \(\chi : X \tto k_u\) that has the property:
  % \begin{quotation}
  for every object \(b\) of \(\mathcal{C}\) and \(\xi : X \tto k_b\) there exists
  one and only one \(g : u \to b\) of \(\mathcal{C}\) that makes
  \[
    \begin{tikzcd}[row sep=tiny]
      u  \ar["g", dd, swap] \\
      & X(i) \ar["{\chi_i}", ul, swap] \ar["{\xi_i}", dl]\\
      b
    \end{tikzcd}
  \]
  commute for every \(i\) in \(\mathcal{I}\).
  % \end{quotation}
\end{definition}

In general, (co)cones can be twisted things and (co)limits, if exist,
can be even more so. Yes, for a functor \(\mathcal{I} \to \mathcal{C}\), the category
\(\mathcal{C}\) has its share, but it is \(\mathcal{I}\) who has the last say in the
research. The role of \(\mathcal{I}\) is to give a \q{shape} of the limits we are
looking for, indeed.

\begin{lemma}\label{lemma:UniqueMorphismToLim}
  Let \(\lambda : k_v \Rightarrow X\) be a limit of a functor
  \(X : \mathcal{I} \to \mathcal{C}\) and \(f, g : a \to v\) in \(\mathcal{C}\). If
  \[
    \lambda_i f = \lambda_i g \quad\text{for every } i \in \obj{\mathcal{I}}
  \]
  then \(f = g\). Dually, if \(\xi : X \Rightarrow k_u\) is a colimit of
  \(X\) and \(f , g : u \to b\) in \(\mathcal{C}\) satisfy
  \[
    f \xi_i = g \xi_i \quad\text{for every } i \in \obj{\mathcal{I}}
  \]
  then \(f = g\).
\end{lemma}

\begin{proof}
  The family of morphisms \(\mu_i := \lambda_i f : a \to X(i)\) is natural in
  \(i\): in fact, for any \(s : i \to j\) in \(\mathcal{I}\) we have
  \(X(s) \mu_i = X(s) \lambda_i f = \lambda_j f = \mu_j\).  By the universal property of
  limit, there exactly one \(a \to v\) that makes the diagrams
  \[
    \begin{tikzcd}[row sep=tiny]
      a \ar["{\mu_i}", dr] \ar["", dd, swap] \\
      & X(i) \\
      v \ar["{\lambda_i}", ur, swap]
    \end{tikzcd}
  \]
  commute. By hypothesis both \(f\) and \(g\) do this, hence they are
  equal.
\end{proof}

\begin{proposition}\label{proposition:LimitsHaveIsomorphicObjects}
  If \(\eta : k_a \Rightarrow X\) and \(\theta : k_b \Rightarrow X\) are limits of the same functor
  \(X : \mathcal{I} \to \mathcal{C}\), then \(a \cong b\). Dually, colimits of the same functor
  have isomorphic vertices.
\end{proposition}

\begin{proof}
  By definition of limit, we a have a unique \(f : a \to b\) and a unique
  \(g : b \to a\) making the triangles in
  \[
    \begin{tikzcd}[row sep=small]
      a \ar["{\eta_i}", dr] \ar["f" description, dd, bend left] \\
      & X(i) \\
      b \ar["{\theta_i}", ur, swap] \ar["g" description, uu, bend left]
    \end{tikzcd}
  \]
  commute for every object \(i\) of \(\mathcal{I}\). Then
  \[
    \begin{aligned}
      & \eta_i = \theta_i f = \eta_i (gf) \\
      & \theta_i = \eta_i g = \theta_i (fg)
    \end{aligned} \quad \text{for every } i \in \obj{\mathcal{I}}
  \]
  Invoking Lemma~\ref{lemma:UniqueMorphismToLim}, we can conclude that
  \(gf = \id_a\) and \(fg = \id_b\).
\end{proof}

In general, one denotes the vertex of any of the limits of a functor
\(X : \mathcal{I} \to \mathcal{C}\) as
\[
  \lim X \quad\text{or}\quad \lim_{i \in \obj{\mathcal{I}}} X(i) .
\]
Dually,
\[
  \colim X \quad\text{or}\quad \underset{i \in \obj{\mathcal{I}}}{\colim} X(i) .
\]
are reserved for the vertex of any of the colimits of \(X\).



\section{Terminal and initial objects}

\begin{definition}[Terminal \& initial objects]
  For \(\mathcal{C}\) category, the limits of the empty functor
  \(\nil \to \mathcal{C}\) are called {\em terminal objects} of \(\mathcal{C}\),
  whereas the colimits {\em initial objects}.  More explicitly:
  \begin{tcbitem}
  \item A {\em terminal object} of \(\mathcal{C}\) is an object \(1\) of
    \(\mathcal{C}\) such that for every object \(x\) of \(\mathcal{C}\) there
    exists one and only one \(x \to 1\) in \(\mathcal{C}\).
  \item An {\em initial object} of \(\mathcal{C}\) is an object \(0\) in
    \(\mathcal{C}\) such that for every object \(x\) in \(\mathcal{C}\) there
    exists one and only one morphism \(0 \to x\) in \(\mathcal{C}\).
  \end{tcbitem}
\end{definition}

\begin{example}[Empty set and singletons]
  It may sound weird, but for every set \(X\) there does exist a
  function \(\nil \to X\); moreover, it is the unique one. To get this,
  think set-theoretically: a function is any subset of \(\nil \times X\) that
  has the property we know. But \(\nil \times X = \nil\), so its unique
  subset is \(\nil\). This set is a function from \(\nil\) to \(X\)
  since the statement
  \begin{quotation}
    for every \(a \in \nil\) there is one and only one \(b \in X\) such that
    \((a, b) \in \nil\)
  \end{quotation}
  is a {\em vacuous truth}. So \(\nil\) is an initial object of
  \(\Set\). This case is quite particular, since the initial objects of
  \(\Set\) are actually equal to \(\nil\). \newline Now let us look for
  terminal objects in \(\Set\). Take an arbitrary set \(X\): there is
  exactly one function from \(X\) to any singleton, that is singletons
  are terminal object of \(\Set\). Conversely, by
  Proposition~\ref{proposition:LimitsHaveIsomorphicObjects}, the
  terminal objects of \(\Set\) must be singletons.
\end{example}

\begin{exercise}
  Trivial groups --- there is a unique way a singleton can be a group ---
  are either terminal and initial objects of \(\Grp\).
\end{exercise}

\begin{example}[The ring \(\integers\)]
  For if \(R\) is a ring, we write \(1_R\) to write the multiplicative
  identity of \(R\). A simple ring homomorphism is
  \[\phi : \integers \to R \,,\ \phi(r) := r 1_R .\]
  Assume now, \(\psi : \integers \to R\) is another ring homomorphism: then
  for every \(r \in \integers\) we have
  \[\psi(r) = r \psi\left(1_R\right) = r \phi\left(1_R\right) = \phi(r) .\]
  Thus we can conclude \(\integers\) is initial in \(\Ring\).
\end{example}

\begin{example}[Recursion]\label{example:Recursion}
  The {\em Recursion Theorem} says:
  \begin{quotation}
    % Let \((\naturals, 0, s)\) be a Peano Model, where
    % \(0 \in \naturals\) and \(s : \naturals \to \naturals\) is its
    % successor function.
    For every set \(X\), \(a \in X\) and \(f : X \to X\) there is a unique
    function \(x : \naturals \to X\) such that
    % \(x_0 = a\) and \(x_{n+1} = f(x_n)\) for every \(n \in \naturals\).
    \begin{align*}
      & x_0 = a \\
      & x_{n+1} = f(x_n) \quad\text{for every } n \in \naturals
    \end{align*}
  \end{quotation}
  % Here, by Peano Model we mean a set \(\naturals\) that has one
  % element, we write \(0\), stood out and a function
  % \(s : \naturals \to \naturals\) such that, all this complying some
  % rules:
  % \begin{tcbenum}
  % \item \(s\) is injective;
  % \item \(s(x) \ne 0\) for every \(x \in \naturals\);
  % \item for if \(A \subseteq \naturals\) has \(0\) and \(s(n) \in A\) for every
  %   \(n \in A\), then \(A = \naturals\).
  % \end{tcbenum}
  We show now how we can involve Category Theory in this case. First of
  all, we need a category where to work. The statement is about things
  made as follows:
  \begin{quotation}
    one set \(X\), one \(x \in X\) and one \(f : X \to X\).
  \end{quotation}
  We may refer to such new things by barely a triple \((X, x, f)\) or by
  \[
    \begin{tikzcd}[column sep=small]
      1 \ar["{x}", r] & X \ar["{f}", r] & X
    \end{tikzcd}
  \]
  where \(1\) is any singleton, as usual.
  % Peano Models are such things, with some additional properties. It is
  % told about the existence and the uniqueness of a certain function.
  We do not want just functions, of course: given
  \[
    % 1 \functo x X \functo f X \text{ and } 1 \functo y Y \functo g Y ,
    \begin{tikzcd}[column sep=small]
      1 \ar["{x}", r] & X \ar["{f}", r] & X
    \end{tikzcd}
    \qquad\text{and}\qquad
    \begin{tikzcd}[column sep=small]
      1 \ar["{y}", r] & Y \ar["{g}", r] & Y
    \end{tikzcd}
  \]
  we take the functions \(r : X \to Y\) such that
  \[
    \begin{tikzcd}[row sep=tiny]
      & X \ar["f", r] \ar["r", dd] & X \ar["r", dd] \\
      1 \ar["x", ur] \ar["y", dr, swap] \\
      & Y \ar["g", r, swap] & Y
    \end{tikzcd}
  \]
  commutes and nothing else. These ones are the things we want to be
  morphisms. Now suppose given
  \[
    \begin{tikzcd}
      & X \ar["f", r] \ar["p", d] & X \ar["p", d] \\
      1 \ar["x", ur] \ar["y", r] \ar["z", dr, swap] & Y \ar["g", r] \ar["q", d] & Y \ar["q", d] \\
      & Z \ar["h", r, swap] & Z
    \end{tikzcd}
  \]
  where all the squares and triangles commute: thus we
  obtain the commuting diagram
  \[
    \begin{tikzcd}[row sep=tiny]
      & X \ar["f", r] \ar["qp", dd] & X \ar["qp", dd] \\
      1 \ar["x", ur] \ar["z", dr, swap] \\
      & Z \ar["h", r, swap] & Z
    \end{tikzcd}
  \]
  This means that composing two morphisms as functions
  in \(\Set\) produces a morphism. This is how we want composition to
  defined in this context. This choice makes the categorial axioms
  automatically respected. We call this category
  \(\mathbf{Peano}\).\newline Being the environment set now, the Recursion Theorem
  becomes more concise:
  \begin{quotation}
    \(\begin{tikzcd}[cramped] 1 \ar["{0}", r] &
      \naturals \ar["{n \mapsto n+1}", r] & \naturals \end{tikzcd}\) is an
    initial object of \(\mathbf{Peano}\).
  \end{quotation}
\end{example}

\begin{exercise}[Induction \(\lrarr\) Recursion]
  In \(\Set\), suppose you have
  \(1 \functo 0 \naturals \functo s \naturals\), where \(s\) is
  injective and \(s(n) \ne 0\) for every \(n \in \naturals\). Demonstrate
  that the following statements are equivalent:
  \begin{tcbenum}
  \item for if \(A \subseteq \naturals\) has \(0\) and \(s(n) \in A\) for every
    \(n \in A\), then \(A = \naturals\);
  \item \(1 \functo 0 \naturals \functo s \naturals\) is an initial
    object of \(\mathbf{Peano}\).
  \end{tcbenum}
\end{exercise}

Limits (colimits) are terminal (initial) objects of appropriate
categories.

\begin{construction}[Category of cones]
  For \(\mathcal{C}\) category, let \(F : \mathcal{I} \to \mathcal{C}\) be a
  functor. Then we define the {\em category of cones} over \(F\) as
  follows.
  \begin{tcbitem}
  \item The objects are the cones over \(F\).
  \item For
    \(\alpha := \set{a \functo{\alpha_i} F(i)}_{i \in \obj{\mathcal{I}}}\) and
    \(\beta := \set{b \functo{\beta_i} F(i)}_{i \in \obj{\mathcal{I}}}\) two cones, the
    morphisms from \(\alpha\) to \(\beta\) are the morphisms
    \(f : a \to b\) of \(\mathcal{C}\) such that
    \[\begin{tikzcd}[row sep=tiny]
        a \ar["{\alpha_i}", dr] \ar["f", dd, swap] \\
        & F(i) \\
        b \ar["{\beta_i}", ur, swap]
      \end{tikzcd}\] commutes for every \(i \in \obj{\mathcal{I}}\).
  \item The composition of morphisms here is the same as that of
    \(\mathcal{C}\).
  \end{tcbitem}
  We write such category as \(\cn_F\). We define also the {\em category
    of cocones} over \(F\), written as \(\cocn_F\).
  \begin{tcbitem}
  \item The objects are the cocones over \(F\).
  \item For
    \(\alpha := \set{F(i) \functo{\alpha_i} a}_{i \in \obj{\mathcal{I}}}\) and
    \(\beta := \set{F(i) \functo{\beta_i} b}_{i \in \obj{\mathcal{I}}}\) cocones, the
    morphisms from \(\alpha\) to \(\beta\) are the morphisms
    \(f : a \to b\) of \(\mathcal{C}\) such that
    \[\begin{tikzcd}[row sep=tiny]
        & a \ar["{\alpha_i}", dl, swap] \ar["f", dd] \\
        F(i) \\
        & b \ar["{\beta_i}", ul]
      \end{tikzcd}\] commutes for every \(i \in \obj{\mathcal{I}}\).
  \item The composition of morphisms here is the same as that of
    \(\mathcal{C}\).
  \end{tcbitem}
  It is quite immediate in either of the cases to show that categorial
  axioms are verified.
\end{construction}

\begin{exercise}
  For \(\mathcal{C}\) category and \(F : \mathcal{I} \to \mathcal{C}\) functor, show that:
  \begin{tcbitem}
  \item limits of \(F\) are terminal objects of \(\cn_F\) and viceversa.
  \item colimits of \(F\) are initial objects of \(\cocn_F\) and
    viceversa.
  \end{tcbitem}
\end{exercise}



\section{Products and coproducts}

\begin{definition}[Products \& coproducts]
  Let \(\mathcal{C}\) be a category, \(I\) a small discrete category (that is
  a set) and consider a functor \(X : I \to \mathcal{C}\). A {\em
    product} of \(X\) is any of the limits of \(X\). Dually, a {\em
    coproduct} of \(X\) is any of the colimits of \(X\). More
  explicitly, recalling that \(X\) in this case
  is just the family \(\set{X_i \mid i \in I}\),
  \begin{itemize}[leftmargin=*]
  \item A product is any family
    \(\set{\proj_i : p \to X_i \mid i \in I}\) of morphisms of
    \(\mathcal{C}\), usually called {\em projections}, with the following
    property: for every family \(\set{f_i : a \to X_i \mid i \in I}\) of
    morphisms of \(\mathcal{C}\) there exists one and only one
    \(h : a \to p\) of \(\mathcal{C}\) such that
    \[
      \begin{tikzcd}[row sep=tiny]
        a \ar["h", dd, swap] \ar["{f_i}", dr] \\
        & X_i \\
        p \ar["{\proj_i}", ur, swap]
      \end{tikzcd}
    \]
    commutes for every \(i \in I\).
  \item A {\em coproduct} of \(\set{X_i \mid i \in I}\) of objects of
    \(\mathcal{C}\) is any family \(\set{\inj_i : X_i \to q \mid i \in I}\) of
    morphisms of \(\mathcal{C}\), often referred to as {\em injections},
    having the property: for every family
    \(\set{g_i : X_i \to b \mid i \in I}\) of morphisms of \(\mathcal{C}\) there
    exists one and only one \(k : q \to b\) of \(\mathcal{C}\) such that
    \[
      \begin{tikzcd}[row sep=tiny]
        b \\
        & X_i \ar["{\inj_i}", dl] \ar["{g_i}", ul, swap] \\
        q \ar["k", uu]
      \end{tikzcd}
    \]
    commutes for every \(i \in I\).
  \end{itemize}
\end{definition}

\begin{example}[Infima and suprema in prosets]
  Consider a proset \((\mathbb P, \le)\) and a subset \(S\) of
  \(\mathbb P\). In this instance, a product of \(S\) is some
  \(p \in \mathbb P\) such that:
  \begin{tcbenum}
  \item \(p \le x\) for every \(x \in S\);
  \item for every \(p' \in \mathbb P\) such that \(p' \le x\) for every
    \(x \in S\) we have \(p' \le p\).
  \end{tcbenum}
  If we have quick look to some existing mathematics, cones over \(S\)
  are what are called {\em lower bounds} of \(S\). An {\em infimum} of
  \(S\) is any of the greatest lower bounds for \(S\).\newline On the other
  hand, a coproduct of \(S\) is some \(q \in \mathbb P\) such that:
  \begin{tcbenum}
  \item \(x \le q\) for every \(x \in S\);
  \item for every \(q' \in \mathbb P\) such that \(x \le q'\) for every
    \(x \in S\) we have \(q \le q'\).
  \end{tcbenum}
  In other words, the cones over \(S\) are precisely the {\em upper
    bounds} of \(S\). A {\em supremum} of \(S\) is any of the lowest
  upper bounds for \(S\).\newline There is a dedicated notations for such
  elements, if \((\mathbb P, \le)\) is a poset: the infimum of \(S\) is
  written as \(\inf S\), whereas \(\sup S\) is the supremum of \(S\). If
  the elements of \(S\) are indexed, that is
  \(S = \set{x_i \mid i \in I}\), then it is customary to write
  \(\inf_{i \in I} x_i\) and \(\sup_{i \in I} x_i\).
\end{example}

\begin{exercise}
  Prosets provide some examples in which some subsets does not have
  infima or suprema.
\end{exercise}

Now let us turn our attention to a pair of quite ubiquitous constructs.

\begin{example}[Cartesian product]\label{example:ProdOfSets}
  Given a family of sets \(\set{X_\alpha \mid \alpha \in \Gamma}\), we have the
  corresponding {\em Cartesian product}
  \[\prod_{\alpha \in \Gamma} X_\alpha := \set{\left. f : \Gamma \to \bigcup_{\alpha \in \Gamma} X_\alpha \right\mid f(\lambda) \in
      X_\lambda \text{ for every } \lambda \in \Gamma} ,\] whose elements are the {\em
    choices} from \(\set{X_\alpha \mid \alpha \in \Gamma}\). As the name indicates, a choice
  \(f\) for every \(\lambda \in \Gamma\) indicates one element of
  \(X_\lambda\). Our product comes with the {\em projections}, one for each
  \(\mu \in \Gamma\),
  \begin{align*}
    & \proj_\mu : \prod_{\alpha \in \Gamma} X_\alpha \to X_\mu \\
    & \proj_\mu(f) := f(\mu) .
  \end{align*}
  Now, any family of functions \(\set{f : A \to X_\alpha \mid \alpha}\) ca be
  compressed into one function
  \[\prod_{\alpha \in \Gamma} f_\alpha : A \to \prod_{\alpha \in I} X_\alpha .\]
  % Taken now any set \(A\) and functions \(f_\alpha : A \to X_\alpha\), one for
  % each \(\alpha \in \Gamma\), we have
  % \[\prod_{\alpha \in \Gamma} f_\alpha : A \to \prod_{\alpha \in \Gamma} X_\alpha\]
  by defining \(\left(\prod_{\alpha \in \Gamma} f_\alpha\right) (a)\) to be the function
  \(\Gamma \to \bigcup_{\alpha \in \Gamma} X_\alpha\) mapping \(\mu \in \Gamma\) to \(f_\mu(a)\).
  % \NotaInterna{Ok, the notation here is becoming cumbersome\dots{}
  % Using \(f_\bullet\) instead of \(\prod_{i \in I} f_i\)? The idea is:
  % \(f_\bullet (a)\) takes \(i\) as input by putting it in place of \(\bullet\).}
  It is simple to show that
  \[\begin{tikzcd}[row sep=small]
      A \ar["{{\prod_{\alpha \in \Gamma}} f_\alpha}", dd, swap] \ar["{f_\mu}", dr] \\
      & X_\mu \\
      \prod_{\alpha \in \Gamma} X_\alpha \ar["{\proj_\mu}", ur, swap]
    \end{tikzcd}\] commutes for every \(\mu \in \Gamma\). Moreover,
  \(\prod_{\alpha \in \Gamma} f_\alpha\) is the only one that does this. Consider any
  function \(g : A \to \prod_{\alpha \in \Gamma} X_\alpha\) with
  \(f_\mu = \proj_\mu g\) for every \(\mu \in \Gamma\): then for every
  \(x \in A\) we have
  \[\begin{aligned}
    \big(g(x)\big)(\mu) &= \proj_\mu \big(g(x)\big) = f_\mu (x) = \\
                      &= p_\mu\left(\left(\prod_{\alpha \in \Gamma} f_\alpha\right)(x)\right) = \left(\left(\prod_{\alpha \in \Gamma} f_\alpha\right)(x)\right)(\mu) ,
  \end{aligned}\]
that is \(g = \prod_{\alpha \in \Gamma} f_\alpha\).
\end{example}

\begin{exercise}
  It may be simple to reason about the Cartesian product of only two
  sets \(X_1\) and \(X_2\). In this case, the product is written as
  \(X_1 \times X_2\) and its elements are represented as pairs
  \((a, b) \in X_1 \times X_2\) rather than functions
  \(f : \set{1, 2} \to X_1 \cup X_2\) with \(f(i) \in X_i\) for
  \(i \in \set{1, 2}\). By setting things like this, the \q{compression}
  of two functions \(f_1 : A \to X_1\) and \(f_2 : A \to X_2\) into a
  function \(A \to X_1 \times X_2\) becomes more obvious.
\end{exercise}

\begin{example}[Coproduct of sets]\label{example:CoprodOfSets}
  For if \(\set{X_\alpha \mid \alpha \in \Lambda}\) is a family of sets, we introduce the
  {\em disjoint union}
  \[\sum_{\alpha \in \Lambda} X_\alpha := \bigcup_{\alpha \in \Lambda} X_\alpha \times \set{\alpha} = \set{(x, \alpha) \mid \alpha \in \Lambda\,,
      \ x \in X_\alpha} .\] While the elements of every member of
  \(\set{X_\alpha \mid \alpha \in \Lambda}\) are amalgamated in
  \(\bigcup_{\alpha \in \Lambda} X_\alpha\), in the disjoint union
  \(\sum_{\alpha \in \Lambda} X_\alpha\) the elements have attached a record of their
  provenience --- in this case, the index of the set they come
  from. Because of this feature, the elements of
  \(\sum_{\alpha \in \Lambda} X_\alpha\) are called {\em dependent pairs}. The disjoint union
  of \(\set{X_\alpha \mid \alpha \in \Lambda}\) has one {\em injection} for each \(\alpha \in \Lambda\):
  \[\inj_\mu : X_\mu \to \sum_{\alpha \in \Lambda} X_\alpha\,, \ \inj_\mu(x) := (x, \mu) .\]
  Similarly to what we have done in the previous example, a family of
  functions \(\set{f_\alpha : X_\alpha \to A \mid \alpha \in \Lambda}\) can be compressed into this
  one
  \begin{align*}
    & \sum_{\alpha \in \Lambda} f_\alpha : \sum_{\alpha \in \Lambda} X_\alpha \to A \\
    & \left(\sum_{\alpha \in \Lambda} f_\alpha\right) (x, \mu) := f_\mu(x) ,
  \end{align*}
  which, in other words, checks the provenience of an element and give
  it to an appropriate function \(f_\alpha\). This new function makes the
  diagram
  \[\begin{tikzcd}[row sep=small]
      A \\
      & X_\mu \ar["{\inj_\mu}", dl] \ar["{f_\mu}", ul, swap] \\
      \sum_{\alpha \in \Lambda} X_\alpha \ar["{\sum_{\alpha \in \Lambda} f_\alpha}", uu]
    \end{tikzcd}\] commute for every \(\mu \in \Lambda\), and it is the unique to
  do this.%\newline
  % Sometimes, you will see written \(\coprod\) instead of \(\sum\), but this
  % does not shift the discourse.
\end{example}

\begin{exercise}
  Prove \(\bigcup_{\alpha \in \Lambda} X_\alpha\) along with appropriate functions
  \[X_\mu \to \bigcup_{\alpha \in \Lambda} X_\alpha\] one for each
  \(\mu \in \Lambda\), is a coproduct if the \(X_\alpha\)-s are pairwise disjoint. If
  some of them are not disjoint, what could go wrong?
  % By the way, \(\bigcup_{\alpha \in \Lambda} X_\alpha\) is isomorphic to an appropriate
  % quotient of \(\sum_{\alpha \in \Lambda} X_\alpha\). Part of the exercise is to find an
  % equivalence relation \(\sim\) on \(\sum_{\alpha \in \Lambda} X_\alpha\) and a function
  % \[\sum_{\alpha \in \Lambda} X_\alpha \to \bigcup_{\alpha \in \Lambda} X_\alpha\]
  % which maps \(\sim\)-equivalent elements to the same element.
\end{exercise}

\begin{exercise}
  Haskell natively offers the function
  \begin{center} {\tt either :: (a -> c) -> (b -> c) -> Either a b -> c}
  \end{center}
  How do {\tt Either a b} and this function fit in the current topic?
  If you accept this little exercise, remember {\tt Either a b} is
  defined to be either {\tt Left a} or {\tt Right b}. We haven't talked
  about the category of types, but it is not be that unseen.
\end{exercise}

\begin{example}[Product of topological spaces]
  Consider now a family of topological spaces \(\set{X_i \mid i \in I}\) and
  let us see if we can have a product of topological space in the sense
  of the Definition above.\newline In order to talk about product topological
  space we shall determine a topology over the set
  \(\prod_{i \in I} X_i\). From the Example~\ref{example:ProdOfSets}, we have
  a nice machinery, but it is all about sets and functions! We define
  the {\em product topology} --- sometimes called \q{Tychonoff topology} ---
  as the smallest among the topologies for \(\prod_{i \in I} X_i\) for which
  all the projections \(\proj_j : \prod_{i \in I} X_i \to X_j\) of the
  Example~\ref{example:ProdOfSets} are continuous.\newline The question is now:
  do these continuous functions form a product in \(\Top\)?  Taking a
  family of continuous functions \(\set{f_i : A \to X_i \mid i \in I}\) and
  looking at the \q{underground} \(\Set\), there does exist one function
  \(\hat f : A \to \prod_{i \in I} X_i\) such that
  \(f_i = \proj_i \hat f\) for every \(i \in I\), but we do not know if it
  is continuous! To give an answer, let us consider the family
  \[\mathcal T := \set{\left. U \subseteq \prod_{i \in I} X_i \text{ open} \right\mid
      \inv{\hat f} U \text{ is open in } A} :\] the idea is that if we
  demonstrate \(\mathcal T\) is a topology for \(\prod_{i \in I} X_i\) and
  \(\mathcal T\) makes all the \(\proj_i\)'s continuous, then we can
  conclude the continuity of \(\hat f\). The first part is immediate, so
  let us focus on the remaining part. If we take an open subset \(V\) of
  \(X_j\), the open subset \(\inv{\proj_j} V\) of the product is in
  \(\mathcal T\), because
  \(\inv{f_j} V = \inv{\hat f} \left(\inv{\proj_j} V\right)\) is open in
  \(A\).
\end{example}

\begin{example}[Coproduct of topological spaces]
  As in the previous example, we move from
  Example~\ref{example:CoprodOfSets}. In Topology, it is maybe more
  customary to use
  \[\coprod_{i \in I} X_i \text{ instead of } \sum_{i \in I} X_i\]
  when \(\set{X_i \mid i \in I}\) is a family of topological spaces. However,
  if we do not give a topology to \(\coprod_{i \in I} X_i\), this object remains
  a bare set. In analogy to what happened with the Cartesian product, we
  prescribe the open subsets of \(\coprod_{i \in I} X_i\) by making reference to
  the injections \(\inj_j : X_j \to \coprod_{i \in I} X_i\):
  \begin{quotation}
    we define a subset \(A\) of \(\coprod_{i \in I} X_i\) to be open if and only
    if \(\inv{\inj_j} A\) is an open subset of \(X_j\) for every
    \(j \in I\).
  \end{quotation}
  Let us recycle the universal property enjoyed by the family of the
  injections, that is for every family
  \(\set{g_i : X_i \to A \mid i \in I}\) of continuous functions there exists
  one function \(\tilde g : \coprod_{i \in I} X_i \to A\) such that
  \(g_j = \tilde g \inj_j\) for every \(j \in I\). Furthermore,
  \(\tilde g\) is continuous: if \(U \subseteq A\) is open, then so are the
  subsets \(\inv{g_j} U \subseteq X_j\); consequently
  \(\inv{g_j} U = \inv{\inj_j} \left(\inv{\tilde g} U\right)\) for every
  \(j \in I\), which implies \(\inv{\tilde g} U\) is open.
\end{example}

\begin{example}[Existential and universal quantifiers]
  Consider one logical predicate \(p\) in one variable, say
  \(x\). Consider a universe of discourse \(\Omega\). \YetToBeTeXed{}
\end{example}

Sometimes, products and coproducts can be isomorphic, as in the
following example.

\begin{example}[Product and coproduct of modules]
  \NotaInterna{Yet to be \TeX{}-ed\dots{}}
\end{example}

\begin{example}[Morphisms out of coproducts of terminal objects]
  A set is an aggregate of single things, as such one thinks about sets
  and as such axioms pertaining sets are crafted. This has a consequence
  on how you do things, and since we are doing Category Theory there is
  something that concerns us:
%
  \begin{quotation}
    How do we introduce a function of sets \(f : A \to B\)?
  \end{quotation}
%
  No surprise in the answer:
%
  \begin{quotation}
    You do it {\em pointwise}: that is for every \(x \in A\) you prescribe
    which element of \(B\) is the image.
  \end{quotation}
%
  Perhaps, you will be more surprised by how Category Theory enters the
  discourse now. As we have said many times, an element \(x \in A\) {\em
    is} the function \(x : 1 \to A\) mapping the unique element of \(1\)
  to \(x\) and applying a function \(f : A \to B\) to an element
  \(x \in A\) means composing \(f x\).
 %
  \begin{quotation}
    The functions \(1 \to A\) form a coproduct in
    \(\Set\). %In other words, if for every function \(x : 1 \to A\) we provide a function \(y_x : 1 \to B\), then there is one and only one \(f : A \to B\) such that \(f x = y_x\).
  \end{quotation}
%
  Make sure you understand why this is the precise formulation of the
  term {\em pointwise} used above. So, here we are:
%
  \begin{quotation}
    Let \(\mathcal{C}\) be a category with a terminal object \(1\) and with
    the property: for every object \(a\) of \(\mathcal{C}\) the elements of
    \(\mathcal{C} (1, a)\) form a coproduct. Then defining morphisms
    \(f : a \to b\) is the same to assigning one morphism \(1 \to b\) for
    every \(1\) the object \(a\) is made of.
  \end{quotation}
  It is the case of \(\Set\) and it is nice, isn't it?
\end{example}

\begin{exercise}
  The previous remark is interesting. Does it work in \(\Top\)? Of
  course no, you are required to think within the categorial frame built
  so far and reinterpret what you already know. You are encouraged to
  explore other categories as well.
\end{exercise}

Let us talk about {\em finite} products, that is products of a finite
set of objects. The following arguments will be useful when we will deal
with finite completeness of categories. Keep an eye on
Figure~\ref{figure:ProductsReductions}.

\begin{figure}
  \centering
  \begin{tabular}{c}
    % \toprule
    reduction from left (Proposition~\ref{proposition:FiniteProdLeftCons}) \\
    % \midrule
    \begin{tikzcd}[ampersand replacement=\&]
      \& \& \& \& \& p_n \ar["{l_n}", dl, swap] \ar["{r_n}", dr] \\
      \& \& \& \& p_{n-1} \ar["{l_{n-1}}", dl, swap] \ar["{r_{n-1}}", dr] \& \& x_n \\
      \& \& \& \cdots{} \ar["{l_4}", dl, swap] \ar["{\cdots{}}", dr] \& \& x_{n-1} \\
      \& \& p_3 \ar["{l_3}", dl, swap] \ar["{r_3}", dr] \& \& \cdots{} \\
      \& p_2 \ar["{l_2}", dl, swap] \ar["{r_2}", dr] \& \& x_3 \\
      x_1 \& \& x_2
    \end{tikzcd}
    \\
    \midrule %\toprule
    reduction from right (Proposition~\ref{proposition:FiniteProdRightCons}) \\
    % \midrule
    \begin{tikzcd}[ampersand replacement=\&]
      \& p_n \ar["{l_n}", dl, swap] \ar["{r_n}", dr] \& \\
      x_1 \& \& p_{n-1} \ar["{l_{n-1}}", dl, swap] \ar["{r_{n-1}}", dr] \\
      \& x_2 \& \& \cdots{} \ar["{\cdots{}}", dl, swap] \ar["{r_4}", dr] \\
      \& \& \cdots{} \& \& p_3 \ar["{l_3}", dl, swap] \ar["{r_3}", dr] \\
      \& \& \& x_{n-2} \& \& p_2 \ar["{l_2}", dl, swap] \ar["{r_2}", dr] \\
      \& \& \& \& x_{n-1} \& \& x_n
    \end{tikzcd} \\
    % \bottomrule
  \end{tabular}
  \caption{Finite products recursively constructed}
  \label{figure:ProductsReductions}
\end{figure}

\begin{proposition}[Finite products, reduction from left]\label{proposition:FiniteProdLeftCons}
  Let \(\mathcal{C}\) be a category a finite set \(\set{x_1, \dots{}, x_n}\),
  with \(n \ge 2\), of objects of \(\mathcal{C}\). Let
  \[\begin{tikzcd}[row sep=tiny, column sep=small]
      & p_2 \ar["{l_2}", dl, swap] \ar["{r_2}", dr] \\
      x_1 & & x_2 \end{tikzcd}\] be one of the products of
  \(\set{x_1, x_2}\) and let
  \[\begin{tikzcd}[row sep=tiny, column sep=small]
      & p_{i+1} \ar["{l_{i+1}}", dl, swap] \ar["{r_{n+1}}", dr] \\
      p_i & & x_{i+1}
    \end{tikzcd}\] be one of the products of \(\set{p_i, x_{i+1}}\).
%
%\[\begin{tikzcd}[sep=small]
%    & & & & & p_n \ar["{l_n}", dl, swap] \ar["{r_n}", dr] \\
%    & & & & p_{n-1} \ar["{l_{n-1}}", dl, swap] \ar["{r_{n-1}}", dr] & & x_n \\
%    & & & \cdots{} \ar["{l_4}", dl, swap] \ar["{\cdots{}}", dr] & & x_{n-1} \\
%    & & p_3 \ar["{l_3}", dl, swap] \ar["{r_3}", dr] & & \cdots{} \\
%    & p_2 \ar["{l_2}", dl, swap] \ar["{r_2}", dr] & & x_3 \\
%    x_1 & & x_2
%  \end{tikzcd}\]
%
  Then the morphisms
  \[\begin{aligned}
    l_2 \cdots{} l_n &: p_n \to x_1 \\
    r_j l_{j+1} \cdots{} l_n &: p_n \to x_j \quad \text{for } j \in \set{2, \dots{}, n-1} \\
    r_n &: p_n \to x_n
  \end{aligned}\]
of \(\mathcal{C}\) do form a product of \(\set{x_1, \dots{}, x_n}\).
\end{proposition}

\begin{proof}
  The proof is conducted by induction on \(n \ge 2\). The case \(n=2\) is
  the base case of our recursive definition. To proceed with the
  inductive step, let us picture the situation:
  \[\begin{tikzcd}
      & & & & p_{n+1} \ar["{l_{n+1}}", dll, swap] \ar["{r_{n+1}}", drr] \\
      & & p_n \ar["{l_2\cdots{}l_n}", dll, swap] \ar["{r_jl_{j+1}\cdots{}l_n}" description, d] \ar["{r_n}", drr] & & & & x_{n+1} \\
      x_1 & & x_j & & x_n \\
      \\
      & & & & & a \ar["{f_1}" description, uulllll] \ar["{f_j}"
      description, uulll] \ar["h" description, uuulll, dotted]
      \ar["{f_n}" description, uul] \ar["k" description, uuuul, dotted]
      \ar["{f_{n+1}}" description, uuur]
    \end{tikzcd}\] where \(j \in \set{2, \dots{}, n-1}\), \(a\) is an
  arbitrary object with morphisms \(f_1, \dots{}, f_n, f_{n+1}\). By the
  universal property of product, we have
  \[\begin{aligned}
    f_1 &= l_1 \cdots{} l_n h \\
    f_j &= r_jl_{j+1} \cdots{} l_n h \\
    f_n &= r_n h
  \end{aligned}\]
for one and only one \(h : a \to p_n\). Again by the universal property of product.
\[\begin{aligned}
  h &= l_{n+1} k \\
  f_{n+1} &= r_{n+1} k
\end{aligned}\]
for a unique \(k : a \to p_{n+1}\). Thus
\[\begin{aligned}
  f_1 &= l_1 \cdots{} l_n l_{n+1} k \\
  f_j &= r_jl_{j+1} \cdots{} l_n l_{n+1} k \\
  f_n &= r_n l_{n+1} k \\
  f_{n+1} &= r_{n+1} k
\end{aligned}\]
and we have concluded.
\end{proof}

\begin{proposition}[Finite products, reduction from right]\label{proposition:FiniteProdRightCons}
  Let \(\mathcal{C}\) be a category a finite set \(\set{x_1, \dots{}, x_n}\),
  with \(n \ge 2\), of objects of \(\mathcal{C}\). Let
  \[\begin{tikzcd}[row sep=tiny, column sep=small]
      & p_2 \ar["{l_2}", dl, swap] \ar["{r_2}", dr] \\
      x_{n-1} & & x_n
    \end{tikzcd}\] be one of the products of \(\set{x_{n-1}, x_n}\) and
  let
  \[\begin{tikzcd}[row sep=tiny, column sep=small]
      & p_{i+1} \ar["{l_{i+1}}", dl, swap] \ar["{r_{n+1}}", dr] \\
      x_{n-i} & & p_i
    \end{tikzcd}\] be one of the products of \(\set{x_{n-i}, p_i}\).
%
%\[\begin{tikzcd}[sep=small]
%    & p_n \ar["{l_n}", dl, swap] \ar["{r_n}", dr] & \\
%    x_1 & & p_{n-1} \ar["{l_{n-1}}", dl, swap] \ar["{r_{n-1}}", dr] \\
%    & x_2 & & \cdots{} \ar["{\cdots{}}", dl, swap] \ar["{r_4}", dr] \\
%%    &     & \cdots{} &          & \cdots{} \\
%    & & \cdots{} & & p_3 \ar["{l_3}", dl, swap] \ar["{r_3}", dr] \\
%    & & & x_{n-2} & & p_2 \ar["{l_2}", dl, swap] \ar["{r_2}", dr] \\
%    & & & & x_{n-1} & & x_n
%  \end{tikzcd}\]
%
  Then the morphisms
  \[\begin{aligned}
    r_2 \cdots{} r_n &: p_n \to x_n \\
    l_j r_{j+1} \cdots{} r_n &: p_n \to x_{n-j+1} \quad \text{for } j \in \set{2, \dots{}, n-1} \\
    r_n &: p_n \to x_1
  \end{aligned}\]
of \(\mathcal{C}\) do form a product of \(\set{x_1, \dots{}, x_n}\).
\end{proposition}

\begin{proof}
  This is \inlinethm{exercise}. You should expect some work like in the
  proof of Proposition~\ref{proposition:FiniteProdLeftCons}.
\end{proof}

\begin{corollary}[Associativity of product]\label{corollary:ProdAssoc}
  In a category \(\mathcal{C}\), let
  \[\begin{tikzcd}[sep=small]
      x_1 & x_1 \times x_2 \ar["{p_1}", l, swap] \ar["{p_2}", r] & x_2
    \end{tikzcd}\] a product of \(\set{x_1, x_2}\),
  \[\begin{tikzcd}[sep=small]
      x_1 \times x_2 & (x_1 \times x_2) \times x_3 \ar["{p_{1,2}}", l, swap]
      \ar["{p_3}", r] & x_3
    \end{tikzcd}\] a product of \(\set{x_1 \times x_2, x_3}\),
  \[\begin{tikzcd}[sep=small]
      x_2 & x_2 \times x_3 \ar["{q_2}", l, swap] \ar["{q_3}", r] & x_3
    \end{tikzcd}\] a product of \(\set{x_2, x_3}\),
  \[\begin{tikzcd}[sep=small]
      x_1 & x_1 \times (x_2 \times x_3) \ar["{q_1}", l, swap] \ar["{q_{2,3}}", r]
      & x_2 \times x_3
    \end{tikzcd}\] a product of \(\set{x_1, x_2 \times x_3}\).
  \[\begin{tikzcd}[sep=small,cramped]
      & & (x_1 \times x_2) \times x_3 \ar["{p_{1,2}}", dl, swap] \ar["{p_3}", ddrr] \\
      & x_1 \times x_2 \ar["{p_1}", dl, swap] \ar["{p_2}", dr] \\
      x_1 & & x_2 & & x_3 \\
      & & & x_2 \times x_3 \ar["{q_2}", ul] \ar["{q_3}", ur, swap] \\
      & & x_1 \times (x_2 \times x_3) \ar["{q_1}", uull] \ar["{q_{2, 3}}", ur,
      swap]
    \end{tikzcd}\] Then
  \[(x_1 \times x_2) \times x_3 \cong x_1 \times (x_2 \times x_3) .\]
\end{corollary}

\begin{proof}
  It follows from Proposition~\ref{proposition:FiniteProdLeftCons} and
  Proposition~\ref{proposition:FiniteProdRightCons}.
\end{proof}

\begin{corollary}\label{corollary:FiniteProductsIffTerminalAndBinaryProducts}
  A category has all finite products if and only if has a terminal
  object and all binary products.
\end{corollary}

\begin{proof}
  One implication is easy. For the opposite one: terminal objects are
  empty products; an object with identity is a product of itself; if you
  are given at least two objects, either of
  Proposition~\ref{proposition:FiniteProdLeftCons} and
  Proposition~\ref{proposition:FiniteProdRightCons} tell you finite
  product are consecutive binary products.
\end{proof}

\begin{exercise}
  In a category with terminal object \(1\), we have
  \(a \times 1 \cong 1 \times a \cong a\).
\end{exercise}



\section{Equalizers and coequalizers}

\begin{definition}[Equalizers \& coequalizers, explicit]
  For \(\mathcal{C}\) category, the limits of functors
  \[
    \left(\begin{tikzcd}[column sep=tiny] \bullet \ar[r, shift left] \ar[r,
        shift right] & \bullet \end{tikzcd}\right) \to \mathcal{C}
  \]
  are called {\em equalizers}; dually, its colimits are called {\em
    coequalizers} instead. More explicitly, given in \(\mathcal{C}\) two
  morphisms
  \begin{equation}
    \begin{tikzcd}
      a \ar["f", r, shift left] \ar["g", r, shift right, swap] & b
    \end{tikzcd}\label{diagram:PairOfParallelMorphisms}
  \end{equation}
  we have:
  \begin{itemize}[leftmargin=*]
  \item An {\em equalizer} of this pair is a morphism \(i : c \to a\)
    such that:
    \begin{enumerate}[leftmargin=*]
    \item The following diagram commutes
      \[
        \begin{tikzcd}
          c \ar["i", r] & a \ar["f", r, shift left] \ar["g", r, shift
          right, swap] & b
        \end{tikzcd}
      \]
    \item For every \(i' : c' \to a\) of \(\mathcal{C}\) making
      \[
        \begin{tikzcd}
          c' \ar["{i'}", r] & a \ar["f", r, shift left] \ar["g", r,
          shift right, swap] & b
        \end{tikzcd}
      \]
      commute, there is one and only one \(k : c' \to c\) in \(\mathcal{C}\)
      such that the following commutes
      \[
        \begin{tikzcd}[row sep=small]
          c' \ar["{i'}", dr] \ar["k", dd, swap] \\
          & a \\
          c \ar["i", ur, swap]
        \end{tikzcd}
      \]
    \end{enumerate}
  \item A {\em coequalizer} of the
    pair~\eqref{diagram:PairOfParallelMorphisms} is a morphism
    \(j : b \to d\) such that
    \begin{enumerate}[leftmargin=*]
    \item The following diagram commutes
      \[
        \begin{tikzcd}
          a \ar["f", r, shift left] \ar["g", r, shift right, swap] & b
          \ar["j", r] & d
        \end{tikzcd}
      \]
    \item For every \(j' : b \to d'\) of \(\mathcal{C}\) making
      \[
        \begin{tikzcd}
          a \ar["f", r, shift left] \ar["g", r, shift right, swap] & b
          \ar["{j'}", r] & d'
        \end{tikzcd}
      \]
      commute, there exists one and only one \(h : d \to d'\) such that
      the following commutes
      \[
        \begin{tikzcd}[row sep=small]
          & d \ar["h", dd] \\
          b \ar["j", ur] \ar["{j'}", dr, swap] \\
          & d'
        \end{tikzcd}
      \]
    \end{enumerate}
  \end{itemize}
\end{definition}

% Before we analyse some example, the following lemma, may be quite
% useful to guide us.



% How could this be of aid? For example, in \(\Set\) this means we have
% to look for inclusions in the domain of the domain of the parallel
% arrows. That is the case, indeed.

The following example explains why the equalizers are called like
that.

\begin{example}[Equalizers in \(\Set\)]
  In \(\Set\), consider two functions
  \[\begin{tikzcd}
      X \ar["f", r, shift left] \ar["g", r, shift right, swap] & Y
    \end{tikzcd} .\] The subset
  \[E := \set{x \in X \mid f(x) = g(x)}\] has the inclusion in \(X\), we
  call it \(i : E \hookrightarrow X\). Of course, we have \(fi = gi\), so one part
  of the work is done. Now, let us take a commuting diagram
  \[\begin{tikzcd}
      E' \ar["{i'}", r] & X \ar["f", r, shift left] \ar["g", r, shift
      right, swap] & Y
    \end{tikzcd} .\] It follows that \(i'(x) \in E\) for every
  \(x \in E'\). Hence, we shall consider the function
  \[h : E' \to E\,, \ h(x) := i'(x) .\] It is immediate now that
  \(i : E \hookrightarrow X\) is an equalizer of \(f\) and \(g\).
\end{example}

Looking at the previous example, the subset where parallel functions
agree gives an equalizer. As a matter of fact, having this intuition in
mind is not wrong.

\begin{lemma}\label{lemma:EqualizersAreMonos}
  Equalizers are monomorphisms. Dually, coequalizers are epimorphisms.
\end{lemma}

\begin{proof}
  Consider
  \[\begin{tikzcd}
      e' \ar["s", r, shift left] \ar["t", r, shift right, swap] & e
      \ar["i", r] & a \ar["f", r, shift left] \ar["g", r, shift right,
      swap] & b
    \end{tikzcd}\] with \(i\) equalizer of \(f\) and \(g\) and
  \(is = it\). We can redraw this diagram as follows:
  \[\begin{tikzcd}[row sep=small]
      e \ar["i", dr] \\
      & a \ar["f", r, shift left] \ar["g", r, shift right, swap] & b \\
      e' \ar["s", uu, bend left=10] \ar["t", uu, bend right=10, swap]
      \ar["{is = it}", ur, swap]
    \end{tikzcd}\] In this case, we have
  \(f(is) = f(it) = g(is) = g(it)\). Thus, by definition of equalizer,
  it must be \(s = t\).
\end{proof}

Thus, if we want one example of coequalizer in \(\Set\), we have to look for
some surjective functions.

\begin{example}[Coequalizers in \(\Set\)]
  In \(\Set\), consider two functions
  \[
    \begin{tikzcd}
      X \ar["f", r, shift left] \ar["g", r, shift right, swap] & Y
    \end{tikzcd} .
  \] Having a commuting diagram like
  \[
    \begin{tikzcd}
      X \ar["f", r, shift left] \ar["g", r, shift right, swap] & Y
      \ar["j", r] & A
    \end{tikzcd}
  \]
  means that we must look for some \(j : Y \to A\) such that, for
  \(x \in X\), the elements \(f(x)\) and \(g(x)\) are brought to the same
  element of \(A\). The thing works if take \(A\) to be the quotient
  \[
    \frac{Y}{f(x) \sim g(x) \ \text{for } x \in X}
  \]
  and define \(j\) as quotient map \(Y \to Y{/}{\sim}\). It only remains to
  verify the universal property of coequalizer. Take any
  \(j' : Y \to Z\) such that \(j'f = j'g\).
  \[
    \begin{tikzcd}[row sep=small]
      & & Y{/}{\sim} \ar[dd, dashed] \\
      X \ar["f", r, shift left] \ar["g", r, shift right, swap] & Y \ar["j", ur] \ar["{j'}", dr, swap] \\
      & & Z
    \end{tikzcd}
  \]
  The existence of the unique function \(Y{/}{\sim} \to Z\) is immediate.
\end{example}

% \begin{exercise}
%   Retrieve Example~\ref{example:PushoutsInSet} and prepare to combine
%   it with the example we have just made. The
%   square~\eqref{diagram:NotYetACommSquareInSet} doesn't even commute,
%   but \(A+B\) with the two injections is a coproduct, not a random
%   thing out there. You can rearrange that diagram too
%   \[\begin{tikzcd}[column sep=large]
%       C \ar["{\linj f}",r, shift left] \ar["{\rinj g}", r, shift
%       right, swap] & A+B
%     \end{tikzcd}\] and summon the canonical projection
%   \(A+B \to \frac{A+B}{\sim}\) which is a coequalizer. At this point, we
%   have the commutative square~\eqref{diagram:NowACommSquareInSet},
%   which we proved to be a pushout square. Luckily, this works in
%   general, and it is up to you to realize why and make the dual of the
%   result too.\newline In a category \(\mathcal{C}\), suppose you have a coproduct
%   \[\begin{tikzcd}a \ar["{\linj}", r] & a+b & b \ar["{\rinj}", l,
%       swap]\end{tikzcd}\] a square
%   \[\begin{tikzcd}
%       c \ar["f", r] \ar["g", d, swap]  & a \ar["\linj", d] \\
%       b \ar["\rinj", r, swap] & a+b
%     \end{tikzcd}\] and a coequalizer \(p : a+b \to d\) of \(\linj f\)
%   and \(\rinj g\). Prove that
%   \[\begin{tikzcd}
%       c \ar["f", r] \ar["g", d, swap]  & a \ar["p\linj", d] \\
%       b \ar["p\rinj", r, swap] & d
%     \end{tikzcd}\] is a pushout square in \(\mathcal{C}\).% Eventually, you have an amazing result.
% \end{exercise}

% \begin{proposition}
%   Categories with binary products/coproducts and
%   equalizers/coequalizers have pullbacks/pushouts.
% \end{proposition}

% Another exercise will help you to prove the converse.

% \begin{exercise}
%   In a category \(\mathcal{C}\), consider two parallel morphisms
%   \[\begin{tikzcd}
%       a \ar["f", r, shift left] \ar["g", r, shift right, swap] & b
%     \end{tikzcd}\] and a product
%   \[\begin{tikzcd}
%       a & a \times b \ar["{p_a}", l, swap] \ar["{p_b}", r] & b
%     \end{tikzcd}\] The universal property of products gives the
%   morphisms \(\bar f, \bar g : a \to a \times b\) such that
%   \(p_a \bar f = \id_a\), \(p_b \bar f = f\), \(p_a \bar g = \id_a\)
%   and \(p_b \bar g = g\):
%   \[\begin{tikzcd}
%       & a \ar["{\id_a}", dl, swap] \ar["{\bar f}", d, shift right, swap] \ar["{\bar g}", d, shift left] \ar["f", dr, shift right, swap] \ar["g", dr, shift left] \\
%       a & a \times b \ar["{p_a}", l] \ar["{p_b}", r, swap] & b
%     \end{tikzcd}\] Consider the pullback square
%   \[\begin{tikzcd}
%       c \ar["m", r] \ar["n", d, swap] & a \ar["{\bar f}", d] \\
%       a \ar["{\bar g}", r, swap] & a \times b
%     \end{tikzcd}\] and prove that:
%   \begin{tcbenum}
%   \item \(m = n\).
%   \item \(m\) is equalizer of \(f\) and \(g\).
%   \end{tcbenum}
% \end{exercise}

% In other words:

% \begin{proposition}
%   Categories with binary products/coproducts have
%   equalizers/coequalizers if and only if have pullbacks/pushouts. That
%   is: in a category with products/coproducts, pullbacks/pushouts can
%   be obtained by equalizers/coequalizers and vice versa.
% \end{proposition}

% \begin{proposition}
%   Categories with binary products/coproducts and pullbacks/pushouts
%   have equalizers/coequalizers.
% \end{proposition}

% \NotaInterna{Write about {\em quotient objects}\dots{}}

% The following proposition may be of help sometimes.

Ok, let us step back to Lemma~\ref{lemma:EqualizersAreMonos}, for
there is something curious. Equalizers are monomorphisms, but what are
the consequences if an equalizer is in addition epic?

\begin{proposition}\label{proposition:EpicEqualizersAreIsomorphisms}
  An epic equalizer is an isomorphism. Dually, a monic epimonomorphism
  is an isomorphism.
\end{proposition}

\begin{proof}
  Assume \(i : e \to a\) is an epic equalizer of
  \[\begin{tikzcd} a \ar["f", r, shift left] \ar["g", r, shift right,
      swap] & b \end{tikzcd}\] From \(fi = gi\), we can derive
  \(f = g\), being \(i\) epic. Hence
  \[\begin{tikzcd} a \ar["{\id_a}", r] & a \ar["f", r, shift left]
      \ar["g", r, shift right, swap] & b \end{tikzcd}\] commutes: by
  the universal property of equalizers, \(\id_a = ik\) for a unique
  \(k : a \to e\). Moreover, by simple computation
  \(iki = i = i \id_e\), which implies that \(ki = \id_e\) since \(i\)
  is monic. In conclusion, \(i\) is invertible.
\end{proof}

How could this become even interesting to us? If a category is such
that every monomorphism is an equalizer of some pair of parallel
arrows, then there monic and epic morphisms are ismorphisms, which we
know it does not occur in every category.

\begin{example}
  In \(\Set\) that phenomenon does occur, let us look in it more
  closely. The problem we have can be stated as follows: take and
  injective function \(f : A \to B\) and find two parallel functions
  parting from \(B\) to which \(f\) is an equalizer. We know, how to
  construct equalizers in \(\Set\) and, even though that is not what
  we want, that example may guide our exploration. The problem we want
  to solve requires to find a certain set \(C\) and two certain
  functions \(h, k : B \to C\). \YetToBeTeXed{}
\end{example}

\begin{exercise}
  In a category \(\mathcal{C}\), if
  \[\begin{tikzcd}
      e \ar["i", r] \ar["i", d, swap] & a \ar["f", d] \\
      a \ar["g", r, swap] & b
    \end{tikzcd}\] is a pullback square, then \(i : e \to a\) is an
  equalizer of the pair \(f\) and \(g\). Is the converse true? If it
  is false, add hypothesis to make it true.
\end{exercise}



\section{Pullbacks and pushouts}

Let \(\mathcal{C}\) be a category. The limits of the functors
\[
  \left(\begin{tikzcd}[sep=tiny]
      \bullet \ar[dr] &         & \bullet \ar[dl] \\
      & \bullet &
    \end{tikzcd}\right) \to \mathcal{C}
\]
are called {\em pullbacks}. Dually, the colimits of the functors
\[
  \left(\begin{tikzcd}[sep=tiny]
      & \bullet \ar[dl] \ar[dr] &         \\
      \bullet & & \bullet
    \end{tikzcd} \right) \to \mathcal{C}
\]
are said {\em pushouts}. More explicitly:

\begin{definition}[Pullbacks \& pushouts, explicit]
  Let \(\mathcal{C}\) a category and a pair of morphisms with the same codomain
  \begin{equation}
    \begin{tikzcd}[row sep=tiny]
      a \ar["f", dr, swap] &   & b \ar["g", dl] \\
      & c &
    \end{tikzcd}
    \label{diagram:MorphismsToPullback}
  \end{equation}
  in \(\mathcal{C}\). A {\em pullback} of~\eqref{diagram:MorphismsToPullback} is
  any pair of morphisms with a common domain
  \[
    \begin{tikzcd}[row sep=tiny]
      & p \ar["h", dl, swap] \ar["k", dr] &   \\
      a & & b
    \end{tikzcd}
  \]
  in \(\mathcal{C}\) such that:
  \begin{tcbitem}
  \item the following square commutes
    \[
      \begin{tikzcd}[row sep=tiny]
        & p \ar["h", dl, swap] \ar["k", dr] &   \\
        a \ar["f", dr, swap] &   & b \ar["g", dl] \\
        & c &
      \end{tikzcd}
    \]
  \item for every
    \[
      \begin{tikzcd}[row sep=tiny]
        & p' \ar["{h'}", dl, swap] \ar["{k'}", dr] &   \\
        a & & b
      \end{tikzcd}
    \]
    in \(\mathcal{C}\) making the diagram
    \[
      \begin{tikzcd}[row sep=tiny]
        & p \ar["{h'}", dl, swap] \ar["{k'}", dr] &   \\
        a \ar["f", dr, swap] &   & b \ar["g", dl] \\
        & c &
      \end{tikzcd}
    \]
    commute there exists one and only one \(r : p' \to p\) in
    \(\mathcal{C}\) such that
    \[
      \begin{tikzcd}
        & p' \ar["{h'}", dl, swap] \ar["r", d, dashed] \ar["{k'}", dr] & \\
        a & p \ar["h", l] \ar["k", r, swap] & b
      \end{tikzcd}
    \]
    commutes.
  \end{tcbitem}
  Assuming now we have a pair of morphisms with the same domain
  \begin{equation}
    \begin{tikzcd}[row sep=tiny]
      & c \ar["f", dl, swap] \ar["g", dr] & \\
      a & & b
    \end{tikzcd}
    \label{diagram:MorphismsToPushout}
  \end{equation}
  in \(\mathcal{C}\), a {\em pushout} of~\eqref{diagram:MorphismsToPushout} is any
  pair of morphisms in \(\mathcal{C}\) with a common codomain
  \[
    \begin{tikzcd}[row sep=tiny]
      a \ar["h", dr, swap] & & b \ar["k", dl] \\
      & q &
    \end{tikzcd}
  \]
  such that:
  \begin{tcbitem}
  \item the square
    \[
      \begin{tikzcd}[row sep=tiny]
        & c \ar["f", dl, swap] \ar["g", dr] &                \\
        a \ar["h", dr, swap] &                                   & b \ar["k", dl] \\
        & q &
      \end{tikzcd}
    \] commutes
  \item for every
    \[
      \begin{tikzcd}[row sep=tiny]
        a \ar["{h'}", dr, swap] & & b \ar["{k'}", dl] \\
        & q' &
      \end{tikzcd}
    \]
    in \(\mathcal{C}\) making
    \[
      \begin{tikzcd}[row sep=tiny]
        & c \ar["f", dl, swap] \ar["g", dr] &                \\
        a \ar["{h'}", dr, swap] &                                   & b \ar["{k'}", dl] \\
        & q &
      \end{tikzcd}
    \]
    commute there exists one and only one \(s : q \to q'\) in
    \(\mathcal{C}\) such that
    \[
      \begin{tikzcd}
        a \ar["h", r] \ar["{h'}", dr, swap] & q \ar["s", d, dashed] & b \ar["k", l, swap] \ar["{k'}", dl] \\
        & q' &
      \end{tikzcd}
    \]
    commutes.
  \end{tcbitem}
\end{definition}

The definitions above become can be more concise though: pullbacks
(pushouts) are products (coproducts) in certain categories.

\begin{proposition}
  A pullback of
  \[\begin{tikzcd}[row sep=tiny]
      a \ar["f", dr, swap] &   & b \ar["g", dl] \\
      & c &
    \end{tikzcd}\] in \(\mathcal{C}\) is any of the products that pair of
  morphisms in \(\mathcal{C} {\downarrow} c\). Dually, a pushout of
  \[\begin{tikzcd}[row sep=tiny]
      & c \ar["f", dl, swap] \ar["g", dr] & \\
      a & & b
    \end{tikzcd}\] in \(\mathcal{C}\) is any of the coproducts such pair of
  morphisms in \(c {\downarrow} \mathcal{C}\).
\end{proposition}

\begin{proof}
  This is \inlinethm{exercise}.
\end{proof}

\begin{exercise}
  Let \(\mathcal{C}\) be a category with initial object \(0\) and terminal
  object \(1\). What are pullbacks of a pair of morphisms with codomain
  \(1\)? What are pushouts of a pair of morphisms with domain \(0\)?
\end{exercise}

\begin{example}[Pullbacks in \(\Set\)]\label{example:PullbacksInSet}
  Now we consider sets and functions as in
  \[
    \begin{tikzcd}
      & A \ar["f", d] \\
      B \ar["g", r, swap] & C
    \end{tikzcd}
  \]
  with the aim to find a pullback of it. From the set
  \[
    D := \set{(a, b) \in A \times B \mid f(a) = g(b)}
  \]
  we can make the functions
  \begin{align*}
    & p : D \to A \\
    & p(a, b) := a \\
    & q : D \to B \\
    & q(a, b) := b
  \end{align*}
  Hence we can draw at least the commuting square
  \[
    \begin{tikzcd}
      D \ar["p", r] \ar["q", d, swap] & A \ar["f", d] \\
      B \ar["g", r, swap] & C
    \end{tikzcd}
  \]
  Now consider
  \[
    \begin{tikzcd}
      D' \ar["{p'}", drr, bend left] \ar["{q'}", ddr, bend right, swap]                 \\
      & D \ar["p", r] \ar["q", d, swap]                                & A \ar["f", d] \\
      & B \ar["g", r, swap] & C
    \end{tikzcd}
  \]
  with \(fp' = gq'\). This implies that
  \(\left(p'(x), q'(x)\right) \in D\) for every \(x \in D'\), and allows us
  to introduce the function
  \begin{align*}
    & r : D' \to D' \\
    & r(x) := \left(p'(x), q'(x)\right)
  \end{align*}
  which is such
  that \(p r = p'\) and \(q r = q'\). Finally, \(r\) is the unique one
  to do so, which fact is immediate for how \(r\) is defined.
\end{example}

Let us remain in \(\Set\) for a moment and explore what happens if we
pull a function back along itself. Provided a function \(f : A \to B\),
the example above tells us we have the pullback square
\[
  \begin{tikzcd}
    R_f \ar["p", r] \ar["q", d, swap] & A \ar["f", d] \\
    A \ar["f", r, swap] & B
  \end{tikzcd}
\]
where
\[
  R_f := \set{(a, b) \in A \times A \mid f(a) = f(b)} .
\]
This subset of \(A \times A\) is an equivalence relation over \(A\), namely
the {\em kernel relation} of \(f\). There is nothing special of \(\Set\)
that prevents us to generalize it to any category \(\mathcal{C}\).

\begin{definition}[Kernel relation]
  In a category \(\mathcal{C}\), the {\em kernel relation} of a morphism
  \(f : a \to b\) is any of the pullbacks of
  \[\begin{tikzcd}
      & a \ar["f", d] \\
      a \ar["f", r, swap] & b
    \end{tikzcd}\]
\end{definition}

In \(\Grp\) pullbacks can be constructed as we have done in \(\Set\). So
what is the kernel relation of a group homomorphism \(f : G \to H\) of \(\Grp\)? As above, it
is the group \(R_f := \set{(a, b) \in G \times G \mid f(a) = f(b)}\) with the
multiplication inherited from \(G \times G\). More explicitly,
\[
  R_f = \left\{ (a, b) \in G \times G \mid a = k b \text{ for some } k \in \ker f
  \right .\}
\]

% As soon as we meet coequalizers, we will have the tool to express the
% quotient \({X}{/}{R_f}\) in a categorial fashion, and thus to motivate
% the general concept of {\em quotient object}.

\begin{example}[Pushouts in \(\Set\)]\label{example:PushoutsInSet}
  Recall what we have done in Example~\ref{example:CoprodOfSets}, but
  change a bit the notation. Take a family of two sets \(A_1\) and
  \(A_2\): write \(A_1+A_1\) instead of using the \(\sum\) or
  \(\coprod\) notation, write \(\linj\) and \(\rinj\) in place of
  \(\inj_1\) and \(\inj_2\), respectively.\newline To get started, let
  us consider sets and functions
  \[\begin{tikzcd}[row sep=tiny]
      & C \ar["f", dl, swap] \ar["g", dr] &   \\
      A & & B
    \end{tikzcd}\] Let us draw a diagram
  \begin{equation}\begin{tikzcd}
      C \ar["f", r] \ar["g", d, swap]  & A \ar["\linj", d] \\
      B \ar["\rinj", r, swap] & A+B
    \end{tikzcd}\label{diagram:NotYetACommSquareInSet}\end{equation}
  By definition of \(A+B\), it can be
  \(\mathtt{left} f(x) \ne \mathtt{right} g(x)\) for some \(x \in
  C\). However, we can make them \q{equal} under an adequate equivalence
  relation \(\sim\): the smallest in which, for \(x \in C\), the
  elements \(\linj f(x)\) and \(\rinj g(x)\) are identified; that is we
  define \(\sim\) to be the smallest equivalence relation containing
  \[R := \set{(\linj f(x), \rinj g(x)) \mid x \in C}.\] In this case,
  let us write \(p\) the projection \(A+B \to \frac{A+B}{\sim}\). The
  new square is
  \begin{equation}\begin{tikzcd}
      C \ar["f", r] \ar["g", d, swap] & A \ar["h", d]    \\
      B \ar["k", r, swap] & \frac{A+B}{\sim}
    \end{tikzcd}\label{diagram:NowACommSquareInSet}\end{equation}
  with \(h := p \linj\) and \(k := p \rinj\) and it is commutative. Now,
  pick
  \[\begin{tikzcd}
      C \ar["f", r] \ar["g", d, swap] & A \ar["h", d] \ar["{h'}", ddr, bend left] \\
      B \ar["k", r, swap] \ar["{k'}", drr, bend right, swap] & \frac{A+B}{\sim} \\
      & & X
    \end{tikzcd}\] such that \(h'f = k'g\). By the universal property of
  coproduct, we have the function
  \[h' + k' : A + B \to X\] such that
  \(\left(h' + k'\right) \linj = h'\) and
  \(\left(h' + k'\right) \rinj = k'\). Taken, for any \(x \in C\), one
  \((\linj f(x), \rinj g(x)) \in R\), we have
  \begin{align*}
    & (h'+g')(\rinj f(x)) = h'f(x) \\
    & (h'+g')(\linj g(x)) = k'g(x) 
  \end{align*}
  which are equal for every \(x \in C\), by
  assumption. \NotaInterna{Talk about generated equivalence relations
    and what follows.} Thus the triangle
  \[\begin{tikzcd}[column sep=tiny]
      A+B \ar["{h'+k'}", rr] \ar["p", dr, swap] & & X \\
      & \frac{A+B}{\sim} \ar[ur, dashed, swap]
    \end{tikzcd}\] commutes for exactly one dashed function. This
  function is the one we are looking for. \NotaInterna{Complete\dots{}}
\end{example}

\begin{exercise}
  In the previous example, what is \(\frac{A+B}{\sim}\) if
  \(C = A \cap B\) and \(f\) and \(g\) are just the inclusions of \(C\)
  in \(A\) and \(B\) respectively? \NotaInterna{There is other material
    to put here\dots{}}
\end{exercise}

\begin{exercise}
  Go back to Example~\ref{example:PushoutsInSet}. What if we started our
  discourse from
  \[\begin{tikzcd}
      C \ar["f", r] \ar["g", d, swap]  & A \ar[d, hookrightarrow] \\
      B \ar[r, hookrightarrow] & A \cup B
    \end{tikzcd}\] instead?
  % You may tackle this question from two different points: trying to
  % readapt the discourse or prove that
  % \((A \cup B){/}{\sim'} \cong (A+B){/}{\sim}\), where \(\sim'\) is
  % the equivalence relation generated by
  % \[f(x) \sim' g(x) \ \text{for } x \in C.\]
\end{exercise}

\begin{lemma}
  Let \(\mathcal{C}\) be a category and
  \begin{equation}\begin{tikzcd}[row sep=tiny]
      a \ar["f", dr, swap] &   & b \ar["g", dl] \\
      & c &
    \end{tikzcd}\label{diagram:MorphismsToBePullbacked}\end{equation}
  a couple of morphisms in \(\mathcal{C}\). If \(\mathcal{C}\) has a product
  \begin{equation}\begin{tikzcd}[row sep=tiny]
      & p \ar["h", dl, swap] \ar["k", dr] &   \\
      a & & b
    \end{tikzcd}\label{diagram:ProductToBePullback}\end{equation}
  of \(a\) and \(b\) such that the square
  \[\begin{tikzcd}[row sep=tiny]
      & p \ar["h", dl, swap] \ar["k", dr] &   \\
      a \ar["f", dr, swap] &   & b \ar["g", dl] \\
      & c &
    \end{tikzcd}\] commutes, then~\eqref{diagram:ProductToBePullback} is
  a pullback of~\eqref{diagram:MorphismsToBePullbacked}.
\end{lemma}

\begin{exercise}[Gluing topological spaces]
  If you are given two spaces \(X\) and \(A\), a subspace
  \(E \subseteq A\) and a continuous function \(f : E \to X\), then
  \(X \sqcup_f A\) denotes the disjoint union \(X \sqcup A\) where every
  \(x \in E\) is identified to \(f(x)\), that
  is% . More formally: on \(X \sqcup A\) we define the equivalence relation \(\sim_f\) as the smallest one containing \(R_f := \set{(f(x), x) \mid x \in E}\) and
  \[X \sqcup_f A := \frac{X \sqcup A}{x \sim f(x) \text{ for } x \in E}
    .\] Find a pushout square
  \[\begin{tikzcd}
      E \ar[r, hookrightarrow] \ar["f", d, swap] & A \ar[d]     \\
      X \ar[r] & X \sqcup_f A
    \end{tikzcd}\] in \(\Top\). The exercise requires you to work about
  the topologies involved and about continuity.
  % with the arrow \(\hookrightarrow\) indicating a mere inclusion. Does
  % this look familiar? Yes! Continue. Part of the exercise is to reason
  % about the topology over the stuff involved here.
\end{exercise}

\begin{example}[CW complexes]
  In Topology, several spaces often employed are homotopic --- or even
  homeomorphic --- to other spaces glued together. Although you can glue
  everything to everything, very simple spaces to attach are disks
  \(\mathbb D^n := \set{x \in \mathbb R^n \mid \lVert x \rVert \le 1}\)
  along their boundaries \(\mathbb S^{n-1} := \partial \dsc^n\). (Pay
  attention to superscripts.) For any topological space \(X\), we can
  perform the following recursive construction:
  \begin{tcbitem}
  \item Let \(X_0\) be the space \(X\) but with the discrete topology.
  \item For \(n \in \naturals\), from topological space \(X_n\) we
    prescribe the construction of another space \(X_{n+1}\). If we are
    given a family \(\set{D_\alpha \mid \alpha \in \Lambda}\) of copies
    of \(\mathbb D^{n+1}\) and collection of continuous functions
    \[\set{\left. f_\alpha : \partial D_\alpha \to X_n \right\mid \alpha
        \in \Lambda}\]
    then we can consider the following topological space
    \[X_{n+1} := \frac{X_n \sqcup \coprod_{\alpha \in \Lambda}
        D_\alpha}{ x \sim f_\alpha(x) \text{ for } \alpha \in \Lambda
        \text{ and } x \in \partial D_\alpha }.\] In other words,
    \(X_{n+1}\) is \(X_n\) with \((n+1)\)-dimensional disks attached to
    it along their boundaries.
  \end{tcbitem}
  As always, we are striving to find some universal property worth of
  consideration. A square comes easily if you consider the inclusions
  running in parallel \(\mathbb S^n \hookrightarrow \mathbb D^{n+1}\)
  and \(X_n \hookrightarrow X_{n+1}\) together with
  \(\mathbb D^{n+1} \hookrightarrow X_{n+1}\) of the construction
  above. The other pieces are the attaching maps: indeed we have a
  commuting square
  % \NotaInterna{A lot of ugly notation and imprecision here\dots{}} As
  % before, at every step of recursive construction just described, we
  % have a commuting square
  \[\begin{tikzcd}[column sep=large]
      \coprod_{\alpha \in \Lambda} \partial D_\alpha \ar[r,
      hookrightarrow] \ar["{\coprod_{\alpha \in \Lambda} f_\alpha}", d,
      swap] & \coprod_{\alpha \in \Lambda} D_\alpha \ar[d,
      hookrightarrow]
      \\
      X_n \ar[r, hookrightarrow] & X_{n+1}
    \end{tikzcd}\] This square is a pushout one in \(\Top\), which is
  easy to prove.
\end{example}

\begin{exercise}[Spheres are CW complexes] Consider the case in which
  \(X_0\) is a single point space and so are the spaces \(X_i\) for
  \(1 \le i \le n-1\). Such situation can be achieved by attaching no
  disk for a while; afterwards, attach one disk \(\dsc^{n}\) along
  \(\sph^{n-1}\) to \(X_{n-1}\). Hence we have a homeomorphism
  \(X_n \cong {\dsc^n}{/}{\sph^{n-1}}\), but you will show something
  more, that is \(X_n \cong \sph^n\).\newline In your Topology course,
  you might have managed to show this as follows:
  \begin{tcbenum}
  \item You have constructed a surjective continuous function
    \(f : \dsc^n \to \sph^n\) and considered the quotient space
    \({\dsc^n}{/}{\sim_f}\) where \(\sim_f\) is the kernel relation
    \NotaInterna{talk about kernel relations!} of \(f\). This relation
    is not a random relation: for every \(x, y \in \dsc^n\), we have
    \(x \sim_f y\) if and only if \(x = y\) or \(x, y \in
    \sph^{n-1}\). As consequence,
    \({\dsc^n}{/}{\sim_f} = {\dsc^n}{/}{\sph^{n-1}}\).
  \item Thanks to the universal property of quotients, the function
    \(\phi : {\dsc^n}{/}{\sph^{n-1}} \to \sph^n\) such that
    \(f = \phi p_n\), with the \(p_n\) the canonical projection, is
    continuous and bijective. Now, recalling that continuous functions
    from compact spaces to Hausdorff spaces are closed, conclude that
    indeed \(\phi\) is a homeomorphism.
  \end{tcbenum}
  The aim of this exercise it that you can arrive to the same result in
  a different manner. If you can recollect your memories or retrieve
  your notes, see if you can recycle the \(f\) above and write a pushout
  square
  \begin{equation}\begin{tikzcd}
      \sph^{n-1} \ar[hookrightarrow, r] \ar["!", d, swap] & \dsc^n \ar["f", d] \\
      X_{n-1} \ar[r] & \sph^n
    \end{tikzcd}\label{diagram:SpherePushout}\end{equation}
  If you do not know how \({\dsc^n}{/}{\sph^{n-1}} \cong \sph^n\), it
  does not matter since you will force yourself to search for a pushout
  square like~\eqref{diagram:SpherePushout}.
\end{exercise}

\begin{exercise}[\(\rpspc^n\) is a CW complex]
  In Topology, the \(n\)-th {\em real projective space} is
  \[\rpspc^n := \frac{\reals^{n+1} \setminus \set{0}}{x \sim \lambda x
      \text{ for } x \in \reals^{n+1}, \lambda \in \reals}\] which is
  known to be homeomorphic to the sphere \(\sph^n\) which has the
  antipodal points identified:
  \[\frac{\sph^n}{\displaystyle x \sim - x \text{ for } x \in \sph^n}
    .\] This observation is the key for the coming arguments. In fact,
  \(\sph^n\) is the boundary of \(\dsc^{n+1}\) and we already have a
  continuous function \(p_n : \sph^n \to \rpspc^n\) that attaches the
  the disc to the projective space along the boundary. Find a pushout
  square of the form
  \[\begin{tikzcd}
      \sph^n \ar[hookrightarrow, r] \ar["p_n", d, swap] & \dsc^{n+1} \ar["?", d] \\
      \rpspc^n \ar[r, hookrightarrow] & \rpspc^{n+1}
    \end{tikzcd}.\]
  % This has remarkable consequence: here is how the real projective
  % spaces are CW complexes!
\end{exercise}

Topology, again, but combined with Group Theory.

\begin{example}[Seifert-van Kampen Theorem]
  Suppose given a topological space \(X\), two open subsets
  \(A, B \subseteq X\) such that \(A \cup B = X\) and one point \(x_0\)
  of \(A \cap B\). Let us denote by \(i_A\), \(i_B\), \(j_A\) and
  \(j_B\) the group morphisms induced by the inclusions
  \(A \cap B \hookrightarrow A\), \(A \cap B \hookrightarrow B\),
  \(A \hookrightarrow X\) and \(B \hookrightarrow X\), respectively. If
  \(A\), \(B\) and \(A \cap B\) are path-connected then,
  \[\begin{tikzcd}[sep=small]
      & \pi_1(A, x_0) \ar["{j_A}", dr] & \\
      \pi_1(A \cap B, x_0) \ar["{i_A}", ur] \ar["{i_B}", dr, swap] & & \pi_1(X, x_0) \\
      & \pi_1(B, x_0) \ar["{j_B}", ur, swap]
    \end{tikzcd}\] is a pushout square of \(\Grp\).
\end{example}

\NotaInterna{The Pullback Lemma is dropped here without a precise plan
  to embed it nicely with examples and further development. It's an
  issue that must be fixed.}

\begin{proposition}[The Pullback Lemma]
  In a category \(\mathcal{C}\) consider a diagram
  \[\begin{tikzcd}
      \bullet \ar["a", r] \ar["d", d, swap] \ar["{Q_1}", dr, phantom] & \bullet \ar["b", r] \ar["g" description, d] \ar["{Q_2}", dr, phantom] & \bullet \ar["c", d]\\
      \bullet \ar["e", r, swap] & \bullet \ar["f", r, swap] & \bullet
    \end{tikzcd}\] where the perimetric rectangle commutes and the
  square on the right is a pullback one. Then that on the left is a
  pullback square is and only if so is the outer rectangle.
\end{proposition}

\begin{proof}
  Let us assume \(Q_1\) is a pullback square first. Consider any choice
  of \(h\) and \(k\) such that \(ch = f(ek)\):
  \[\begin{tikzcd}
      \bullet \ar["h", bend left=10, drrr, blue] \ar["k", bend right=10, ddr, swap, blue] & & & \\
      & \bullet \ar["a", r, swap] \ar["d", d] \ar["{Q_1}", dr, phantom] & \bullet \ar["b", r, swap] \ar["g", d] \ar["{Q_2}", dr, phantom] & \bullet \ar["c", d, blue]\\
      & \bullet \ar["e", r, swap, blue] & \bullet \ar["f", r, swap,
      blue] & \bullet
    \end{tikzcd}\] Being \(Q_2\) a pullback square, there exists one and
  only one \(l\) such that \(h = bl\) and \(gl = ek\).
  \[\begin{tikzcd}
      \bullet \ar["h", bend left=10, drrr, blue] \ar["k", bend right=10, ddr, swap, blue] \ar["l"{description}, bend left=8, drr] & & & \\
      & \bullet \ar["a", r, swap] \ar["d", d] \ar["{Q_1}", dr, phantom] & \bullet \ar["b", r, swap] \ar["g", d] \ar["{Q_2}", dr, phantom] & \bullet \ar["c", d, blue]\\
      & \bullet \ar["e", r, swap, blue] & \bullet \ar["f", r, swap,
      blue] & \bullet
    \end{tikzcd}\] We have just said that this square in red commutes:
  \[\begin{tikzcd}
      \bullet \ar["h", bend left=10, drrr] \ar["k", bend right=10, ddr, swap, red] \ar["l"{description}, drr, bend left=8, red] & & & \\
      & \bullet \ar["a", r, swap] \ar["d", d] \ar["{Q_1}", dr, phantom] & \bullet \ar["b", r, swap] \ar["g", d, red] \ar["{Q_2}", dr, phantom] & \bullet \ar["c", d]\\
      & \bullet \ar["e", r, swap, red] & \bullet \ar["f", r, swap] &
      \bullet
    \end{tikzcd}\] Now, being \(Q_1\) a pullback square, we have one
  \(m\) such that\(l = am\) and \(k = dm\):
  \[\begin{tikzcd}
      \bullet \ar["h", bend left=10, drrr] \ar["k", bend right=10, ddr, swap, red] \ar["l"{description}, drr, bend left=8, red] \ar["m"{description}, dr] & & & \\
      & \bullet \ar["a", r, swap] \ar["d", d] \ar["{Q_1}", dr, phantom] & \bullet \ar["b", r, swap] \ar["g", d, red] \ar["{Q_2}", dr, phantom] & \bullet \ar["c", d]\\
      & \bullet \ar["e", r, swap, red] & \bullet \ar["f", r, swap] &
      \bullet
    \end{tikzcd}\] At this point, we have \(b a m = b l = h\) and
  \(dm = k\). To conclude the first half of the theorem, you have to
  pick any \(m'\) making commute the triangles in green:
  \[\begin{tikzcd}
      \bullet \ar["h", bend left=10, drrr, green!65!black] \ar["k", bend right=10, ddr, swap, green!65!black] \ar["am" description, drr, bend left=8] \ar["{m'}" description, dr, green!65!black] & & & \\
      & \bullet \ar["a", r, swap, green!65!black] \ar["d", d, green!65!black] \ar["{Q_1}", dr, phantom] & \bullet \ar["b", r, swap, green!65!black] \ar["g", d] \ar["{Q_2}", dr, phantom] & \bullet \ar["c", d]\\
      & \bullet \ar["e", r, swap] & \bullet \ar["f", r, swap] & \bullet
    \end{tikzcd}\] Being \(Q_2\) a pullback square, we have
  \(am' = am\). In conclusion, being \(Q_1\) a pullback square too, we
  have \(m = m'\). \NotaInterna{Finer explanation here\dots{}}
\end{proof}

\begin{exercise}
  Prove the remaining part of the theorem above.
\end{exercise}

\begin{example}[Character functions]
  Consider a subset \(A\) of some larger set \(X\). You sure know a
  simple function called {\em character function} with just says if an
  element of \(X\) is a member of \(A\):
  \[\chi_A : X \to \set{\mathtt{true}, \mathtt{false}}\,,\ \chi_A (x)
    := \begin{cases} \mathtt{true} & \text{if } x \in A \\
      \mathtt{false} & \text{otherwise} \end{cases}\] From now on, let
  us write \(\Omega\) to mean \(\set{\mathtt{true},
  \mathtt{false}}\). As always, let us draw what we have:
\[\begin{tikzcd}
    A \ar[r, hookrightarrow] & X \ar["{\chi_A}", d] \\
    & \Omega
  \end{tikzcd}\] with \(A \hookrightarrow X\) being the usual
inclusion. The composition of such functions is function constant to
\(\mathtt{true}\). We know that constant functions are such because they
can be factored through some function \(A \to 1\) \NotaInterna{write
  about this explicitly somewhere}, which results in a commuting square
\[\begin{tikzcd}
    A \ar[r, hookrightarrow] \ar["!", d, swap] & X \ar["{\chi_A}", d] \\
    1 \ar["{\lambda x. \mathtt{true}}", r, swap] & \Omega
  \end{tikzcd}\] Well, this square is a pullback square. Of course, that
is not all we have to say. \NotaInterna{To be continued.}
\end{example}

% Let us introduce a small generalization of pullbacks and pushouts. Let
% \(\mathcal{I}\) be a category which has one object \(b\) and, for
% \(\lambda \in \Gamma\), one object \(x_\lambda\) one morphism
% \(f_\lambda : x_\lambda \to b\); there is no other morphisms than these ones and the
% identities in \(\mathcal{I}\). What are limits of functors
% \(F : \mathcal{I} \to \mathcal{C}\) in more explicit terms? Any limit of
% \(F\) consists of one object \(p\), one morphism \(h : p \to b\) and, for
% every \(\lambda \in \Gamma\), one morphism \(g_\lambda : p \to x_\lambda\) such that
% \[f_\lambda g_\lambda = h \text{ for every } \lambda \in \Gamma .\] We
% shall call this kind of limits {\em generalized
%   pullbacks}. \NotaInterna{\q{Unofficial} name.}  Generalized because
% the usual definition of pullback is obtained by choosing \(\Gamma\) to
% be a set of solely two elements. Observe also, as in the case of
% pullbacks, generalized pullbacks are products in certain comma
% categories. \NotaInterna{Talk about generalized pushouts too.}

% \begin{example}[Generalized pullbacks in \(\Set\)]
%   Consider a family of sets \(\set{X_\lambda \mid \lambda \in \Gamma}\), one set
%   \(T\) and functions \(f_\lambda : X_\lambda \to T\).
%   \[E := \set{\left. x \in \prod_{\lambda \in \Gamma} X_\lambda \right\mid f_\alpha (x(\alpha)) = f_\beta (x(\beta))
%       \text{ for every } \alpha, \beta \in \Gamma} .\] We have then commutative squares
%   \[\begin{tikzcd}
%       E \ar["{p_\alpha}", r] \ar["{p_\beta}", d, swap] & X_\alpha \ar["{f_\alpha}", d] \\
%       X_\beta \ar["{f_\beta}", r, swap] & T
%     \end{tikzcd}\] where the \(p_\alpha\)-s are restrictions to \(E\) of the
%   \(\proj_\alpha\)-s in Example~\ref{example:ProdOfSets}. Consider, for
%   \(\alpha, \beta \in \Gamma\),
%   \[\begin{tikzcd}
%       Y \ar["{g_\alpha}", drr, bend left] \ar["{g_\beta}", ddr, bend right, swap] \\
%       & E \ar["{p_\alpha}", r] \ar["{p_\beta}", d, swap] & X_\alpha \ar["{f_\alpha}", d] \\
%       & X_\beta \ar["{f_\beta}", r, swap] & T
%     \end{tikzcd}\] where \(f_\alpha g_\alpha = f_\beta g_\beta\). Thus
%   \(g_\bullet (y) \in E\) for every \(y \in Y\), which fact motivates the
%   following function
%   \[Y \to E\,, \ y \to g_\bullet(y).\] After realizing how here we have
%   generalized Example~\ref{example:PullbacksInSet}, continue and finish
%   this example: \inlinethm{exercise}.
% \end{example}

\begin{example}[Generalized pushouts in \(\Set\)]
  Recall Example~\ref{example:PushoutsInSet}, because we will need
  it. Consider a family of sets \(\set{X_\mu \mid \mu \in \Lambda}\),
  one set \(T\) and functions \(g_\mu : S \to X_\mu\), one for each
  \(\mu \in \Lambda\). As in the binary case, we have non commuting
  diagrams
  \[\begin{tikzcd}
      S \ar["{g_\delta}", r] \ar["{g_\eta}", d, swap] & X_\eta \ar["{\inj_\delta}", d] \\
      X_\eta \ar["{\inj_\eta}", r, swap] & {\displaystyle\sum_{\mu \in
          \Lambda} X_\mu}
    \end{tikzcd}\] for \(\delta, \eta \in \Lambda\). This is just to
  push ourselves to the next move: consider the smallest relation
  \(\sim\) on the disjoint sum containing
  \[R := \set{(\inj_\delta g_\delta (x), \inj_\eta g_\eta(x)) \mid x \in
      S \text{ and } \delta, \eta \in \Lambda} .\] In this case, we have
  a commutative diagram
  \[\begin{tikzcd}
      S \ar["{g_\delta}", r] \ar["{g_\eta}", d, swap] & X_\delta \ar["{q_\delta}", d] \\
      X_\eta \ar["{q_\eta}", r, swap] & \frac{\sum_{\mu \in \Lambda}
        X_\mu}{\sim}
    \end{tikzcd}\] where, for \(\mu \in \Lambda\), we define
  \(q_\mu (x)\) to be the \(\sim\)-equivalence class of \(x \in
  X_\mu\). Take now one function \(h_\mu : X_k \to F\), for
  \(\mu \in \Lambda\), such that \(h_\delta g_\delta = h_\eta g_\theta\)
  for every \(\delta, \eta \in \Lambda\).
  \[\begin{tikzcd}
      S \ar["{g_\delta}", r] \ar["{g_\eta}", d, swap] & X_\delta \ar["{q_\delta}", d] \ar["{h_\delta}", ddr, bend left] \\
      X_\eta \ar["{q_\eta}", r, swap] \ar["{h_\eta}", drr, bend right, swap] & \frac{\sum_{\mu \in \Lambda} X_\mu}{\sim} \ar[dr, dashed] \\
      & & F
    \end{tikzcd}\] To construct the dotted function, we proceed
  similarly as in the binary case. We have the function
  \[\sum_{\mu \in \Lambda} h_\mu : \sum_{\mu \in \Lambda} X_\mu \to F\]
  which satisfies
  \begin{align*}
    & \left( \sum_{\mu \in \Lambda} h_\mu \right) (\inj_\delta g_\delta (x)) = h_\delta g_\delta (x) =  %\\
    & = h_\eta g_\eta (x) = \left( \sum_{\mu \in \Lambda} h_\mu \right) (\inj_\eta g_\eta (x))
  \end{align*}
  \[\left( \sum_{\mu \in \Lambda} h_\mu \right) (\inj_\delta g_\delta (x)) = h_\delta g_\delta (x) = h_\eta
    g_\eta (x) = \left( \sum_{\mu \in \Lambda} h_\mu \right) (\inj_\eta g_\eta (x))\] for
  \(x \in S\) and \(\delta, \eta \in \Lambda\). The function we are looking for is the one
  induce by \(\sum_{\mu \in \Lambda} h_\mu\). The continuation of this example is
  \inlinethm{exercise}.
\end{example}

\begin{exercise}
  Can you do something similar to what we have done with finite
  products? We are referring to
  Proposition~\ref{proposition:FiniteProdLeftCons},
  Proposition~\ref{proposition:FiniteProdRightCons} and
  Corollary~\ref{corollary:ProdAssoc}.
\end{exercise}



\section{(Co)Completeness}

As the examples of the previous sections have shown, existence of limits
in general is not to be taken for granted.

\begin{definition}
  A category \(\mathcal{C}\) is said to be {\em (co)complete} whenever any functor
  \(X : \mathcal{I} \to \mathcal{C}\) has a (co)limit. Dually,
  \(\mathcal{C}\) is said {\em finitely (co)complete} when every functor
  \(\mathcal{I} \to \mathcal{C}\) out of a {\em finite} category \(\mathcal{I}\) admits a (co)limit.
\end{definition}

Let us look at the first non trivial case, \(\Set\). Let \(X : \mathcal{I} \to \mathcal{C}\)
be any functor. Write \(\left\lvert X \right\rvert : \obj{\mathcal{I}} \to \mathcal{C}\) the
{\em underlying discrete functor}, that is the functor defined on objects by \(\left\lvert X
\right\rvert (a) :=  X(a)\).

Why are we introducing this construction? Well, look at the definitions
of limit and product: they are almost the same. The cones of the
definition of limits are just families of morphisms \(\pi_i : P \to X(i)\) for \(i
\in \mathcal{I}\). In contrast, the cones of the definition of limit are families of
morphisms \(\lambda_i : L \to X(i)\) for \(i \in \mathcal{I}\) such that \(X(f) \lambda_i = \lambda_j\)
for every \(i, j \in \mathcal{I}\) and \(f : i \to j\) in \(\mathcal{I}\). The case
\(\mathcal{C} = \Set\) can be illuminating in this sense. 

We know \(\Set\) has products, that is \(\left\lvert X \right\rvert\)
has a limit: the set
\[
  \prod_{i \in \mathcal{I}} \left\lvert X \right\rvert (i) = \prod_{i \in \mathcal{I}} X(i) 
\]
with the family of functions
\begin{align*}
  & \pi_k : \prod_{i \in \mathcal{I}} X(i) \to X(k) \\
  & \pi_k (x) := x_k
\end{align*}
Let us see if within \(\prod_{i \in \mathcal{I}} X(i)\) we can find a limit of the
diagram \(X : \mathcal{I} \to \mathcal{C}\). Well, since we want naturality, let us just
consider
\begin{align*}
  L &:= \left\{ \left. x \in \prod_{i \in \mathcal{I}} X(i) \right\mid X(f) \pi_i (x) = \pi_j (x)
      \text{ for every } i, j \in \mathcal{I} \text{ and } f \in \mathcal{I}(i, j) \right\} = \\
    &= \left\{ \left. x \in \prod_{i \in \mathcal{I}} X(i) \right\mid X(f) (x_i) = x_j
      \text{ for every } i, j \in \mathcal{I} \text{ and } f \in \mathcal{I}(i, j) \right\}
\end{align*}
Then
\begin{align*}
  & \lambda_i : L \to X(i) \\
  & \lambda_i(x) := x_i
\end{align*}
is natural in \(i\) and gives a limit of \(X\).

\begin{proposition}[Completeness Theorem]\label{proposition:Completeness}
  Categories that have products and equalizers are complete.
\end{proposition}

\begin{proof}
  Take \({\bf I}\) and \(\mathcal{C}\) to be two categories, and
  \(X : {\bf I} \to \mathcal{C}\) any functor; just for convenience, let us
  write \(I\) for \(\obj{{\bf I}}\). \(\mathcal{C}\) has all products, so
  let us write
  \[\set{\left. p \functo{\phi_i} X_i \right\mid i \in I}\]
  for one of --- it does not matter which one, right? --- the products of
  \(\set{X_i \mid i \in I}\). The class
  \[H := \set{j \in I \mid {\bf I}(i, j) \ne \nil \text{ for some } i \in I}\]
  will be useful for the constructions to come. For \(i \in I\),
  \(j \in H\) and \(f \in {\bf I}(i, j)\) we can draw this
  \[\begin{tikzcd}[row sep=tiny]
      & X_i \ar["{X_f}", dr, bend left=15pt] \\
      p \ar["{\phi_i}", ur, bend left=15pt] \ar["{\phi_j}", rr, swap, bend
      right] & & X_j
    \end{tikzcd}\] Again by the fact that \(\mathcal{C}\) has products, let
  us write
  \[\set{\left. q \functo{\xi_j} X_j \right\mid j \in H}\]
  for one of the products of \(\set{X_j \mid j \in H}\). Now, the universal
  property of products yields two morphisms
  \begin{equation}
    \begin{tikzcd} p \ar["\delta", r, shift left] \ar["\theta", r, shift right, swap] & q \end{tikzcd}
    \label{diagram:TwoProducts}
  \end{equation}
  of \(\mathcal{C}\) obtained as follows:
  \begin{tcbenum}%[label=(\arabic*), ref=\arabic*]
  \item\label{universal:theta} \(\theta\) is the one that factors
    \(\phi_j\) through \(\xi_j\) for \(j \in H\), viz \(\phi_j = \xi_j \theta\).
  \item\label{universal:delta} \(\delta\) is the unique that factors
    \(X_f \phi_i\) through \(\xi_j\) for every \(i \in I\),
    \(j \in H\) and \(f \in {\bf I}(i, j)\), that is \(X_f \phi_i = \xi_j \delta\)
  \end{tcbenum}
  \(\mathcal{C}\) has equalizers too, so let \(\epsilon : e \to p\) be one of the
  equalizers of the parallel morphisms
  in~\eqref{diagram:TwoProducts}. Now that everything is arranged, the
  rest of the proof is to prove that
  \[\set{\left. e \functo{\phi_i \epsilon} X_i \right\mid i \in I}\]
  is a limit of \(X\). It is important, however, to check
  preliminarily that it is a natural transformation. Take
  \[\begin{tikzcd}[row sep=tiny]
      & X_i \ar["{X_f}", dd] \\
      e \ar["{\phi_i\epsilon}", ur] \ar["{\phi_j\epsilon}", dr, swap] \\
      & X_j
    \end{tikzcd}\] with \(i, j \in I\) and \(f \in {\bf I}(i, j)\). If
  \(j \notin H\), then the commutativity of the diagram is a vacuous truth;
  otherwise,
  \[X_i \phi_i \epsilon = \underbrace{\xi_j \delta \epsilon = \xi_j \theta \epsilon}_{\epsilon \text{ is
        equalizer}} = \phi_j \epsilon .\] So, let us conclude the proof:
  provided a natural transformation
  \[\set{\left. a \functo{\sigma_i} X_i \right\mid i \in I} ,\]
  we show how to construct \(a \to e\) that makes
  \[\begin{tikzcd}[row sep=tiny]
      a \ar["{\sigma_i}", dr] \ar[dd] \\
      & X_i \\
      e \ar["{\phi_i \epsilon}", ur, swap]
    \end{tikzcd}\] commute. By universal property of product, there is
  a unique \(\mu : a \to p\) such that
  \(\sigma_i = \phi_i \mu = \xi_i \theta \mu\) for every \(i \in I\). In particular, for
  \(j \in H\), \(i \in I\) and \(f \in {\bf I} (i, j)\)
  \[
    \sigma_j =
    \begin{cases}
      \phi_j \mu = \xi_j \theta \mu & \text{by~\eqref{universal:theta}} \\
      X_f \sigma_i = X_f \phi _i \mu = \xi_j  \delta \mu & \parbox{9.5em}{because \(\sigma\) is
                                        a natural transformation
                                        and~\eqref{universal:delta}}
    \end{cases}
  \]
  As a consequence of the universal property of product of \(\xi\), we
  must have \(\theta \mu = \delta \mu\). Moreover, being
  \(\epsilon : e \to p\) an equalizer of~\eqref{diagram:TwoProducts}, then
  \(\mu = \epsilon \psi\) for exactly one \(\psi : a \to e\) of
  \(\mathcal{C}\). Thus \(\sigma_i = \phi_i \mu = \phi_i \epsilon \psi\), so
  \(\psi\) is what we are are looking for; at this point you can observe
  the uniqueness of \(\psi\) as well. That's all.
\end{proof}

% \begin{lemma}
%   A category has finite products if and only if it has a terminal
%   object and all binary products.
% \end{lemma}

% \begin{proof}
%   One implication is trivial. Let \(\mathcal{C}\) be a category that has a
%   terminal object, say \(1\), and all binary products
%   \NotaInterna{write what this would mean}. Let
%   \(\set{x_1, \dots{}, x_n} \subseteq \obj{\mathcal{C}}\) and construct one product
%   for them. We proceed by induction on \(n\). If \(n=0\), we have a
%   product: the terminal object \(1\). Assume now
%   \(\set{x_1, \dots{}, x_n}\) has a product, say the set of morphisms
%   of \(\mathcal{C}\)
%   \[\set{\alpha_i : p \to x_i \mid i = 1, \dots{}, n} .\]
%   Take one \(x_{n+1} \in \obj{\mathcal{C}}\). By assumption, there is a
%   product of \(p\) and \(x_{n+1}\), a pair of morphisms
%   \[p \xleftarrow \beta q \xrightarrow \gamma x_{n+1} .\] The question is: do
%   the morphisms \(\alpha_i \beta : q \to x_i\), for
%   \(i = 1, \dots{}, n\), and \(\gamma\) form a product of
%   \(\set{x_1, \dots{}, x_n, x_{n+1}}\)? Yes, \inlinethm{exercise} (the
%   drawing below is a hint).
%   \[\begin{tikzcd}
%       & & r \ar["{f_i}", ddll, swap, bend right] \ar["{f_{n+1}}", ddr,
%       bend left]
%       \ar[ddl, bend right, dashed] \ar[d, dotted] \\
%       & & q \ar["\beta", dl] \ar["\gamma", dr, swap] \\
%       x_i & p \ar["{\alpha_i}", l] & & x_{n+1}
%     \end{tikzcd}\qedhere\]
% \end{proof}

A special place is for finite (co)limits.

\begin{proposition}[Finite Completeness Theorem I]
  Categories having terminal objects, binary products and equalizers
  are finitely complete. \NotaInterna{Write a definition for \q{finite
      completeness}.}
\end{proposition}

\begin{proof}
  Use
  Corollary~\ref{corollary:FiniteProductsIffTerminalAndBinaryProducts}
  and the argument to prove the Completeness Theorem.
\end{proof}

% Actually, we have another finite completeness theorem, which requires
% some preliminary work.

% \begin{lemma}
%   If a category has a terminal object and pullbacks, then it has
%   binary products and equalizers.
% \end{lemma}

% \begin{proof}
%   Call \(\mathcal{C}\) the category of the assumptions and \(1\) one of its
%   terminal objects. An exercise in the section about pullbacks asked
%   you to show that the pullbacks of \(a \to 1 \gets b\) are products of
%   \(a\) and \(b\). Now we show how to get an equalizer out of a
%   terminal object and an appropriate pullback. Consider two parallel
%   morphisms
%   \[\begin{tikzcd}
%       a \ar["f", r, shift left] \ar["g", r, shift right, swap] & b
%     \end{tikzcd}\] of \(\mathcal{C}\). Now, let
%   \[\begin{tikzcd}a & a \times b \ar["{p_a}", l, swap] \ar["{p_b}", r] &
%       b \end{tikzcd}\] be one the products of \(a\) and
%   \(b\). Afterwards, define \(\bar f : a \to a \times b\) to be that morphism
%   such that \(p_a \bar f = \id_a\) and \(f = p_b \bar f\); similarly,
%   let \(\bar g : a \to a \times b\) be the morphism such that
%   \(p_a \bar g = \id_a\) and \(g = p_b \bar g\).
%   \[\begin{tikzcd}
%       & a \ar["{\id_a}", dl, swap] \ar["{\bar f}", d] \ar["f", dr] \\
%       a & a \times b \ar["{p_a}", l] \ar["{p_b}", r, swap] & b
%     \end{tikzcd} \quad \begin{tikzcd}
%       & a \ar["{\id_a}", dl, swap] \ar["{\bar g}", d] \ar["g", dr] \\
%       a & a \times b \ar["{p_a}", l] \ar["{p_b}", r, swap] & b
%     \end{tikzcd}\] By assumption, \(\mathcal{C}\) has the pullback square
%   \[\begin{tikzcd}
%       p \ar["m", r] \ar["n", d, swap] & a \ar["{\bar f}", d] \\
%       a \ar["{\bar g}", r, swap] & a \times b
%     \end{tikzcd}\] Here \(m = \alpha \bar f m = \alpha \bar g n = n\), so we can
%   collapse all to the commuting
%   \[\begin{tikzcd}
%       p \ar["m", r] & a \ar["{\bar f}", r, shift left] \ar["{\bar g}",
%       r, shift right, swap] & a \times b
%     \end{tikzcd}\] with \(m\) equalizer of \(\bar f\) and \(\bar g\).
% \end{proof}

% \begin{exercise}
%   Show that \(m\) is an equalizer. (Hint: \(\bar f\) and \(\bar g\)
%   are monomorphisms.)
% \end{exercise}

\begin{proposition}[Finite Completeness Theorem II]
  Categories that have terminal objects and pullbacks are finitely
  complete.
\end{proposition}

\begin{proof}
  Use the previous Lemma and the Finite Completeness Theorem I.
\end{proof}

Let us sum all up in one corollary:

\begin{corollary}[Finite Completeness Theorem]
  For any category, the following facts are equivalent:
  \begin{tcbenum}
  \item it is finitely complete
  \item it has a terminal, binary products and equalizers
  \item it has a terminal object and pullbacks
  \end{tcbenum}
\end{corollary}



\section{Exercises}

In the category \(\Modu_R\), homsets have the structure of abelian
group. We want to do generalise this.

\begin{definition}[Preadditve categories]
  A {\em preadditive category} \(\mathcal{C}\) is a category where
  \begin{enumerate}[leftmargin=*]
  \item For every \(x, y \in \obj{\mathcal{C}}\) there are an operation
    \(+_{x, y}\) on \(\mathcal{C}(x, y)\) and a {\em zero morphism}
    \(0_x^y \in \mathcal{C}(x, y)\) making \(\mathcal{C}(x, y)\) an abelian group. In general,
    we write \(+\) without pedices when there is no ambiguity to
    unclutter our notation. Sometimes, we do the same for zero
    morphisms.
  \item For every \(a, b, c \in \obj{\mathcal{C}}\) and \(f : a \to b\) the functions
    \begin{align*}
      & f_\ast := \mathcal{C} (c, f) : \mathcal{C}(c, a) \to \mathcal{C}(c, b) \\
      & f^\ast := \mathcal{C} (f, c) : \mathcal{C}(b, c) \to \mathcal{C}(a, c)
    \end{align*}
    are group homomorphisms.
  \end{enumerate}
\end{definition}

In the following exercise you will prove elementary properties about
limits and colimits in preadditive categories.

\begin{exercise}
  Let \(\mathcal{C}\) be a preadditive category.
  \begin{enumerate}
  \item In \(\mathcal{C}\), initial objects are also terminal; dually, terminal
    objects are initial. Thus initial objects and terminal objects are
    {\em zero objects}. We usually write zero objects in preadditive
    categories as \(0\). Pay attention, when \(0\) denotes zero
    morphisms.

  \item If \(\mathcal{C}\) has zero object \(0\), then
    \(0_a^b : a \to b\) is the composition of
    \[
      \begin{tikzcd}[column sep=small]
        a \ar["{\exists!}", r] & 0 \ar["{\exists!}", r] & b
      \end{tikzcd}
    \]

  \item In \(\mathcal{C}\) let
    \[
      \begin{tikzcd}[column sep=small]
        a & a \times b \ar["{p_a}", l, swap] \ar["{p_b}", r] & b
      \end{tikzcd}
    \]
    be a product. Then introduce two morphisms
    \begin{align*}
      & i_a := \langle\id_a, 0_a^b\rangle : a \to a \times b \\
      & i_b := \langle0_b^a, \id_b\rangle : a \to a \times b
    \end{align*}
    Show these two morphisms form a coproduct. Hence in a preaddtive
    category \(a \times b \cong a + b\). Now do the converse: take a coproduct
    and find a way to have a product. It is customary to write \(a \oplus b\)
    for both \(a \times b\) and \(a + c\).

  \item In \(\mathcal{C}\) consider objects and morphisms
    \[
      \begin{tikzcd}[column sep=small]
        a \ar["{i_a}", r, swap, shift right] & c \ar["{p_a}", l,
        swap, shift right] \ar["{p_b}", r, shift left] & b \ar["{i_b}",
        l, shift left]
      \end{tikzcd}
    \]
    such that
    \begin{align*}
      & p_a i_b = 0_b^a \\
      & p_b i_a = 0_a^b \\
      & i_a p_a + i_b p_b = \id_c
    \end{align*}
    Show that \(p_a\) and \(p_b\) are a product and \(i_a\) and \(i_b\)
    are a coproduct.
  \end{enumerate}
\end{exercise}

\begin{exercise}[Limits in \(\Set^{\mathcal{C}}\)] %\([\mathcal{C}, \Set]\)]
  In this exercise you will understand how to compute the limits of
  diagrams \(X : \mathcal{I} \to [\mathcal{C}, \Set]\). Let us unbox
  things and see where to put hands first. We need a functor \(L :
  \mathcal{C} \to \Set\) together with a family of natural
  transformations \(\lambda_i : L \Rightarrow X(i)\) natural in \(i \in \obj{\mathcal{I}}\). (Make
  sure you got it right: for every \(i\) we have a natural
  transformation \(\lambda_i : L \Rightarrow X(i)\) and \(\lambda_i\) is natural in
  \(i\). Write things in components.) Also, this cone \(\lambda\) should enjoy a
  universal property: for every cone \(\left\{ \eta_i : A \Rightarrow X(i) \mid i \in
    \obj{\mathcal{I}} \right\}\) there is a unique \(\psi : A \Rightarrow L\) such that \(\lambda_i \psi
  = \eta_i\) for every \(i\). If you write this in components, you would
  have the commutative diagrams
  \[
    \begin{tikzcd}[row sep=small]
      A(a) \ar["{\psi_a}", dd, swap] \ar["{\eta_{i,a}}", dr] \\
      & X(i)(a) \\
      L(a) \ar["{\lambda_{i, a}}", ur, swap]
    \end{tikzcd}
  \]
  So what if \(L(a) := \lim_{i \in \obj{\mathcal{I}}} X(i)(a)\)? You will be ask to
  not hurry too much and to duly verify all the details.
  \begin{enumerate}
  \item For every \(a \in \obj{\mathcal{C}}\) construct a functor \(X(-)(a) :
    \mathcal{C} \to \Set\) defined on objects by \(X(-)(a) (i) := X(i)(a)\). These
    functors have a limit. Why?
  \item For every \(a \in \obj{\mathcal{C}}\) let
    \[
      \left\{ \lambda_{i, a} : L(a) \to X(i)(a) \mid i \in \obj{\mathcal{I}} \right\}
    \]
    a limit for \(X(-)(a)\). Can you find a functor \(L : \mathcal{C} \to \Set\)
    such that \(L(a)\) is a the vertex of the limit cone above?
  \item Finally, verify that \(\left\{ \lambda_i : L \Rightarrow X(i) \mid i \in \obj{\mathcal{I}}
    \right\}\) is a limit of \(X : \mathcal{I} \to [\mathcal{C}, \Set]\).
  \end{enumerate}
\end{exercise}


%%% Local Variables:
%%% mode: LaTeX
%%% TeX-master: "../CT"
%%% TeX-engine: luatex
%%% End:
